% ============================================================
%  算術的宇宙の物理学 第2版
%  ― Spec(Z) からスカラー・テンソル重力まで
%  Wright Brothers Research Program 公開教科書
%  wb_textbook_v2.tex (master file)
% ============================================================
\documentclass[11pt,a4paper]{report}

% ---- Fonts (XeLaTeX) ----
\usepackage{fontspec}
\usepackage{xeCJK}
\setCJKmainfont{Hiragino Mincho ProN}
\setCJKsansfont{Hiragino Sans}

% ---- Mathematics ----
\usepackage{amsmath,amssymb,amsthm,mathtools,array}

% ---- Layout ----
\usepackage[margin=25mm]{geometry}
\usepackage{setspace}
\onehalfspacing

% ---- Tables & Lists ----
\usepackage{booktabs}
\usepackage{longtable}
\usepackage{enumitem}

% ---- Graphics & Color ----
\usepackage{xcolor}
\usepackage{tcolorbox}
\tcbuselibrary{skins,breakable}
\usepackage{tikz}
\usetikzlibrary{arrows.meta,calc,positioning,decorations.pathmorphing}

% ---- Hyperlinks ----
\usepackage[unicode,pdfencoding=auto,
  colorlinks=true,linkcolor=blue!60!black,
  citecolor=green!50!black,urlcolor=purple!60!black]{hyperref}
\pdfstringdefDisableCommands{%
  \def\cdots{...}%
  \def\Z{Z}%
  \def\Q{Q}%
  \def\R{R}%
  \def\Spec{Spec}%
  \def\Tr{Tr}%
  \def\cd{cd}%
  \def\Gal{Gal}%
  \def\KMS{KMS}%
  \def\Mpl{M_Pl}%
  \def\Geff{G_eff}%
  \def\Lw{L_w}%
  \def\\{}%
  \def\quad{ }%
  \def\textbf#1{#1}%
}

% ---- Headers ----
\usepackage{fancyhdr}
\pagestyle{fancy}
\fancyhf{}
\fancyhead[L]{\small\sffamily Wright Brothers テキストブック 第2版}
\fancyhead[R]{\small\sffamily 算術的宇宙の物理学}
\fancyfoot[C]{\thepage}

% ============================================================
%  Theorem environments (日本語)
% ============================================================
\newtheorem{theorem}{定理}[chapter]
\newtheorem{proposition}[theorem]{命題}
\newtheorem{corollary}[theorem]{系}
\newtheorem{lemma}[theorem]{補題}
\newtheorem{definition}[theorem]{定義}
\theoremstyle{remark}
\newtheorem{remark}[theorem]{注意}
\newtheorem{example}[theorem]{例}
\renewcommand{\proofname}{証明}

% ============================================================
%  tcolorbox styles
% ============================================================

% 重要結果（緑）
\newtcolorbox{keyresult}[1][]{%
  colback=green!5,colframe=green!60!black,
  fonttitle=\bfseries,title={#1},breakable}

% 定義ボックス（青）
\newtcolorbox{dfnbox}[1][]{%
  colback=blue!5,colframe=blue!50!black,
  fonttitle=\bfseries,title={#1},breakable}

% 警告（赤）
\newtcolorbox{warningbox}[1][]{%
  colback=red!5,colframe=red!50!black,
  fonttitle=\bfseries,title={#1},breakable}

% 実験（黄）
\newtcolorbox{experimentbox}[1][]{%
  colback=yellow!5,colframe=yellow!60!black,
  fonttitle=\bfseries,title={#1},breakable}

% 高校生向け補足（オレンジ）
\newtcolorbox{hsbox}[1][高校生ノート]{%
  colback=orange!5,colframe=orange!50!black,
  fonttitle=\bfseries,title={#1},breakable}

% 仮説ボックス（紫）
\newtcolorbox{hypbox}[1][]{%
  colback=purple!5,colframe=purple!50!black,
  fonttitle=\bfseries,title={#1},breakable}

% ステータス（青）
\newtcolorbox{statusbox}[1][]{%
  colback=blue!5,colframe=blue!50!black,
  fonttitle=\bfseries,title={#1},breakable}

% Go/Stop（赤）
\newtcolorbox{gostopbox}[1][]{%
  colback=red!5,colframe=red!60!black,
  fonttitle=\bfseries,title={#1},breakable}

% ============================================================
%  Custom commands
% ============================================================
\newcommand{\Z}{\mathbb{Z}}
\newcommand{\Q}{\mathbb{Q}}
\newcommand{\R}{\mathbb{R}}
\newcommand{\C}{\mathbb{C}}
\newcommand{\N}{\mathbb{N}}
\newcommand{\Spec}{\operatorname{Spec}}
\newcommand{\Tr}{\operatorname{Tr}}
\newcommand{\cd}{\operatorname{cd}}
\newcommand{\Gal}{\operatorname{Gal}}
\newcommand{\KMS}{\mathrm{KMS}}
\newcommand{\Mpl}{M_{\mathrm{Pl}}}
\newcommand{\Geff}{G_{\mathrm{eff}}}
\newcommand{\Lw}{\mathcal{L}_{\mathrm{w}}}

% ---- Misc ----
\setcounter{tocdepth}{2}

% ============================================================
\begin{document}

% ---- Title page ----
\begin{titlepage}
\centering
\vspace*{2cm}
{\Huge\bfseries 算術的宇宙の物理学}\\[0.8cm]
{\Huge\bfseries 第2版}\\[1.2cm]
{\LARGE Spec($\Z$) からスカラー・テンソル重力まで}\\[2.5cm]
{\Large Wright Brothers Research Program 公開教科書}\\[2cm]
{\large 2026年4月}\\[3cm]
{\normalsize GitHub: \texttt{isshii-jpeg/wright-brothers}}\\[0.5cm]
{\small 本書はオープンアクセスです。自由に配布・引用できます。}
\end{titlepage}

\tableofcontents
\newpage

% ============================================================
%  PART I: 算術的時空の3公理
% ============================================================
\part{算術的時空の3公理}
% ============================================================
%  Part I: 算術的時空の3公理
%  wb_textbook_v2_axioms.tex
%  Chapters 1--3
% ============================================================

% ############################################################
\chapter{Spec(\texorpdfstring{$\Z$}{Z}) --- 整数のスペクトル}
\label{ch:specZ}
% ############################################################

この章では、本書全体の出発点となる数学的対象 $\Spec(\Z)$ を導入する。
高校数学の「素数」から出発し、代数幾何学が素数に与える幾何学的構造を解説する。
最終目標は、\textbf{整数の環 $\Z$ が3次元的な幾何を持つ}という驚くべき事実（Artin--Verdier の定理）
に到達することである。


% ============================================================
\section{素数とは何か}
\label{sec:primes-intro}
% ============================================================

\begin{hsbox}[高校生ノート：素数の復習]
\textbf{素数}（prime number）とは、$1$ より大きい自然数で、
$1$ と自分自身以外に正の約数を持たない数のことである。
最初のいくつかを列挙しよう：
\[
  2,\; 3,\; 5,\; 7,\; 11,\; 13,\; 17,\; 19,\; 23,\; 29,\; \ldots
\]
\textbf{算術の基本定理}（Fundamental Theorem of Arithmetic）により、
すべての自然数 $n \ge 2$ は素数の積として\emph{一意的に}（順序を除いて）表される：
\[
  n = p_1^{a_1} p_2^{a_2} \cdots p_k^{a_k}.
\]
例えば $360 = 2^3 \cdot 3^2 \cdot 5$ である。
素数は自然数の「原子」---これ以上分解できない基本構成要素---と考えてよい。
\end{hsbox}

素数が無限に存在することは、紀元前3世紀にユークリッドが証明した。
その論法は美しく単純である：素数が有限個 $p_1, \ldots, p_k$ しかないと仮定すると、
$N = p_1 p_2 \cdots p_k + 1$ はどの $p_i$ でも割り切れない。
しかし $N \ge 2$ なので素因数を持つはずであり、矛盾する。

\begin{remark}
この証明の構造---\emph{全体を掛け合わせて $1$ を足す}---は、
後に登場する $\zeta$ 関数のオイラー積表示と深い関係がある。
\end{remark}

素数の分布に関する最も基本的な問は「$x$ 以下の素数はいくつあるか」である。
素数計数関数 $\pi(x) = \#\{p \le x : p \text{ は素数}\}$ に対して、
\textbf{素数定理}（Prime Number Theorem）は
\[
  \pi(x) \sim \frac{x}{\ln x} \quad (x \to \infty)
\]
を主張する。より精密な情報はリーマンゼータ関数の零点分布に符号化されている。


% ============================================================
\section{オイラー積とリーマンゼータ関数}
\label{sec:euler-product}
% ============================================================

\begin{definition}[リーマンゼータ関数]
\label{def:zeta}
$\Re(s) > 1$ に対して、\textbf{リーマンゼータ関数}を
\[
  \zeta(s) = \sum_{n=1}^{\infty} \frac{1}{n^s}
\]
と定義する。
\end{definition}

この級数は $\Re(s) > 1$ で絶対収束する。
$\zeta$ 関数の最も重要な性質は、算術の基本定理の解析的反映である\textbf{オイラー積}表示である。

\begin{theorem}[オイラー積, 1737]
\label{thm:euler-product}
$\Re(s) > 1$ において
\[
  \zeta(s) = \prod_{p:\text{素数}} \frac{1}{1 - p^{-s}}.
\]
\end{theorem}

\begin{proof}
各素数 $p$ に対して等比級数
\[
  \frac{1}{1 - p^{-s}} = 1 + p^{-s} + p^{-2s} + p^{-3s} + \cdots
\]
を用いる。有限個の素数 $p_1, \ldots, p_k$ について積をとると、
算術の基本定理により、$p_1, \ldots, p_k$ のみを素因数に持つすべての自然数 $n$ にわたる和
$\sum_{n} n^{-s}$ が得られる。$k \to \infty$ とすれば $\zeta(s) = \sum_{n=1}^{\infty} n^{-s}$
に収束する。
\end{proof}

\begin{hsbox}[高校生ノート：オイラー積の意味]
オイラー積は「掛け算の世界（素因数分解）と足し算の世界（級数）を結ぶ橋」である。
$\zeta$ 関数の一つ一つの因子 $(1-p^{-s})^{-1}$ が素数 $p$ に対応し、
それらの積が\emph{すべての}自然数にわたる和と等しくなる。
これが本書の物理的解釈の出発点となる：
\begin{center}
\emph{各素数 $p$ は独立な自由度であり、\\
全体の分配関数はそれらの積で書ける。}
\end{center}
\end{hsbox}

$\zeta(s)$ は $\Re(s) > 1$ の領域から $\C$ 全体に解析接続され、
$s = 1$ にのみ単純極を持つ。\textbf{リーマン予想}は、
自明でない零点がすべて $\Re(s) = 1/2$ 上にあると主張する。
本書では直接リーマン予想を必要としないが、ゼータ関数の零点は
後に（Part IIで）暗黒エネルギー密度 $\Omega_\Lambda$ の導出に現れる。


% ============================================================
\section{$\Spec(\Z)$ の幾何学}
\label{sec:specZ-geometry}
% ============================================================

代数幾何学では、環 $R$ に対してその\textbf{素イデアルのスペクトル}
\[
  \Spec(R) = \{\mathfrak{p} \subset R : \mathfrak{p} \text{ は素イデアル}\}
\]
を考え、これをザリスキ位相により位相空間とする。

\begin{definition}[$\Spec(\Z)$]
\label{def:specZ}
$\Z$ の素イデアルは以下の通りである：
\begin{itemize}
  \item $(0)$：\textbf{生成的点}（generic point）。$\Z$ の商体 $\Q$ に対応する。
  \item $(p)$（$p$ は素数）：\textbf{閉点}。有限体 $\mathbb{F}_p = \Z/p\Z$ に対応する。
\end{itemize}
従って $\Spec(\Z)$ は集合として
\[
  \Spec(\Z) = \{(0)\} \cup \{(2), (3), (5), (7), (11), \ldots\}
\]
であり、閉点の集合は素数の集合と一対一に対応する。
\end{definition}

\begin{hsbox}[高校生ノート：「空間」としての整数]
通常、「空間」と聞くと $\R^3$ のような連続的な広がりを想像する。
しかし $\Spec(\Z)$ は素数という離散的な「点」から構成される空間である。
にもかかわらず、この空間には豊かな幾何学的構造がある。

比喩的に言えば：
\begin{itemize}
  \item 素数 $p$ は空間上の「点」
  \item 生成的点 $(0)$ はすべての点に「染み渡る」---どの閉点の近傍にも含まれる
  \item $\Spec(\Z)$ 全体は一つの「線」のように振る舞う
\end{itemize}
しかし、次節で見るように、エタール・コホモロジーの意味ではこの「線」は
3次元的な構造を持つ。
\end{hsbox}

$\Spec(\Z)$ 上の構造層 $\mathcal{O}_{\Spec(\Z)}$ は、
各開集合 $U$ に対して「$U$ 上の正則関数の環」を割り当てる。
例えば、点 $(p)$ における茎（stalk）は局所化 $\Z_{(p)}$、
生成的点 $(0)$ における茎は $\Q$ である。

\begin{remark}
代数幾何学では $\Spec(\Z)$ は\textbf{最終対象}（final object）の役割を果たす。
すなわち、すべてのスキーム $X$ に対して一意的な射 $X \to \Spec(\Z)$ が存在する。
この意味で $\Spec(\Z)$ は「すべての幾何学の母体」であり、
物理学で「時空の最も基本的な基盤」と解釈する動機はここにある。
\end{remark}


% ============================================================
\section{Artin--Verdier の定理：$\cd(\Spec(\Z)) = 3$}
\label{sec:artin-verdier}
% ============================================================

$\Spec(\Z)$ のエタール・コホモロジー的次元は古典的な結果として知られている。

\begin{theorem}[Artin--Verdier]
\label{thm:artin-verdier}
$\Spec(\Z)$ のエタール・コホモロジー的次元は $3$ である：
\[
  \cd(\Spec(\Z)) = 3.
\]
すなわち、すべてのエタール層 $\mathcal{F}$ に対して $H^n_{\text{ét}}(\Spec(\Z), \mathcal{F}) = 0$
（$n > 3$）であり、$n = 3$ でゼロにならない層が存在する。
\end{theorem}

この定理は深い意味を持つ。$\Spec(\Z)$ はナイーブには1次元的対象（「線」）に見えるが、
エタール位相の意味では\textbf{3次元多様体}のように振る舞う。

\begin{keyresult}[物理的解釈：時空の次元]
Artin--Verdier の定理 $\cd(\Spec(\Z)) = 3$ は、
$\Spec(\Z)$ 上のコホモロジーが3次元閉多様体と同じ構造を持つことを意味する。
本書のプログラムでは、これを時空の\textbf{空間次元が $3$ である理由}
の算術的起源と解釈する：
\[
  \boxed{d_{\text{空間}} = \cd(\Spec(\Z)) = 3.}
\]
このアイデアは第7章以降のスカラー・テンソル重力理論の基盤となる。
\end{keyresult}

Artin--Verdier の定理のもう一つの顕著な帰結は、
$\Spec(\Z)$ 上のエタール・コホモロジーが\textbf{ポアンカレ双対性}に類似した
双対定理を満たすことである。これは $\Spec(\Z)$ が3次元閉多様体の
算術的類似物であるという視点を補強する。


% ============================================================
\section{素数と結び目のアナロジー}
\label{sec:primes-knots}
% ============================================================

\begin{hsbox}[高校生ノート：素数が結び目のように振る舞う]
$\Spec(\Z)$ が3次元多様体のように振る舞うという事実から、
\textbf{算術位相幾何学}（arithmetic topology）と呼ばれる分野が生まれた。
この分野の核心にあるのは以下の辞書である：

\begin{center}
\begin{tabular}{lll}
\toprule
\textbf{数論} & \textbf{3次元位相幾何学} & \textbf{対応} \\
\midrule
$\Spec(\Z)$ & 3次元多様体 $M^3$ & 全体空間 \\
素数 $(p)$ & 結び目 $K \subset M^3$ & 基本的対象 \\
$\Spec(\mathbb{F}_p)$ & 円 $S^1$ & 閉点／芯 \\
ガロア群 & 基本群 $\pi_1$ & 対称性 \\
相互律 & 絡み数 & 「絡み合い」の測定 \\
\bottomrule
\end{tabular}
\end{center}

素数 $p$ は3次元空間の中の「結び目」（knot）に対応する！
結び目同士が「絡み合う」（linking）ように、素数同士も二次相互律などを通じて
互いに「絡み合う」のである。

この比喩は単なるアナロジーではなく、エタール基本群・被覆空間の
レベルで数学的に精密化できる。
\end{hsbox}

Mazur \cite{Mazur1964}、Morishita \cite{Morishita2012} らの仕事により、
この辞書は極めて豊かな構造を持つことが知られている。
本書の物理的プログラムにおいて、この辞書は次の問いを自然にする：

\begin{quote}
\emph{もし $\Spec(\Z)$ が3次元多様体ならば、\\
その上の「物理」---場の理論、重力、ゲージ理論---はどのようなものか？}
\end{quote}

この問いに答えるのが本書の目標であり、
次章ではそのための出発点となる3つの公理を定式化する。

\bigskip

\begin{remark}[本章のまとめ]
\begin{enumerate}
  \item 素数は自然数の「原子」であり、ゼータ関数のオイラー積を通じて解析学に接続する。
  \item $\Spec(\Z)$ は素数を「点」とする幾何学的空間である。
  \item Artin--Verdier の定理により $\cd(\Spec(\Z)) = 3$、
        すなわちこの空間は3次元多様体のように振る舞う。
  \item 算術位相幾何学の辞書は「素数 $\leftrightarrow$ 結び目」の対応を与え、
        $\Spec(\Z)$ 上の物理を考える動機を与える。
\end{enumerate}
\end{remark}



% ############################################################
\chapter{3つの公理から $\zeta(t)$ へ}
\label{ch:three-axioms}
% ############################################################

本章は第2版の\textbf{最も重要な新内容}である。
前章で導入した $\Spec(\Z)$ の幾何学を出発点として、
\textbf{たった3つの物理的公理}からリーマンゼータ関数 $\zeta(t)$
がプランクスケールのヒートカーネルとして一意的に導出されることを示す。

この導出は本書のプログラム全体の\textbf{論理的基盤}であり、
以降の章で現れるすべての物理的帰結---暗黒エネルギー密度、微細構造定数、
質量比、スカラー・テンソル重力---はこの3公理から流れ出す。


% ============================================================
\section{動機：プランクスケールの自由度}
\label{sec:axioms-motivation}
% ============================================================

現代物理学では、プランク長 $\ell_{\mathrm{Pl}} \sim 10^{-35}\,\mathrm{m}$
以下のスケールにおける時空の構造は不明である。
量子重力のあらゆるアプローチ---弦理論、ループ量子重力、
因果集合理論、非可換幾何学---は、
プランクスケールで時空が何らかの\textbf{離散的構造}を持つことを示唆する。

本書のアプローチは以下の問いから始まる：

\begin{quote}
\emph{プランクスケールの離散的自由度が自然数 $n \in \N$ でラベルされるとき、\\
最小限の仮定から物理がどこまで導出できるか？}
\end{quote}

この問いに対する答えを与えるのが、以下の3つの公理である。


% ============================================================
\section{公理1：算術離散性}
\label{sec:axiom1}
% ============================================================

\begin{dfnbox}[公理1（算術離散性）]
プランクスケールの自由度は自然数 $n \in \N = \{1, 2, 3, \ldots\}$ でラベルされる。
系のヒートカーネル（分配関数的対象）は、重み関数 $w : \N \to \C$ を用いて
\begin{equation}
  K(t) = \sum_{n=1}^{\infty} w(n)\, n^{-t}
  \label{eq:heat-kernel-general}
\end{equation}
と書ける。ここで $t > 0$ は（逆温度に相当する）パラメータであり、
各モード $n$ のエネルギーは $E_n = \log n$ で与えられる。
\end{dfnbox}

\paragraph{公理1の意味.}
この公理は2つのことを主張している：
\begin{enumerate}
  \item プランクスケールのモードは\textbf{離散的}であり、自然数で番号付けられる。
  \item 各モード $n$ のエネルギーは $\log n$ である（$n^{-t} = e^{-t \log n}$）。
\end{enumerate}

エネルギーが $E_n = \log n$ であるという点は一見奇妙に見えるかもしれないが、
自然数 $n$ の「大きさ」を測る最も自然な尺度が対数であることを思い出せば、
これは $n$ を「プランクエネルギーの $\log n$ 倍」と読むことに対応する。

\begin{hsbox}[高校生ノート：ヒートカーネルとは？]
ヒートカーネル（heat kernel）は、もともと「熱がどう伝わるか」を記述する関数である。
金属棒の各点の温度の時間発展は熱方程式 $\partial_t u = \Delta u$ で記述され、
その解は初期条件とヒートカーネルの畳み込みで書ける。

量子力学や場の理論では、ヒートカーネル $K(t) = \Tr(e^{-tH})$ は
ハミルトニアン $H$ のスペクトル情報をすべて含む「指紋」のようなものである。
$t$ が小さいときの $K(t)$ の振る舞いから、
空間の体積・曲率・位相的情報が読み取れる（Seeley--DeWitt 展開、第3章で詳述）。

公理1は、この $K(t)$ がディリクレ級数 $\sum w(n) n^{-t}$ の形をしていると仮定する。
\end{hsbox}


% ============================================================
\section{公理2：素数独立性}
\label{sec:axiom2}
% ============================================================

\begin{dfnbox}[公理2（素数独立性）]
重み関数 $w$ は\textbf{乗法的}（multiplicative）である：
$\gcd(m, n) = 1$ ならば
\begin{equation}
  w(mn) = w(m)\, w(n).
  \label{eq:multiplicativity}
\end{equation}
\end{dfnbox}

\paragraph{公理2の意味.}
乗法性は「$m$ と $n$ が互いに素ならば、モード $mn$ の重みは $m$ と $n$ の
重みの積である」と述べている。算術の基本定理により、すべての自然数は
$n = p_1^{a_1} \cdots p_k^{a_k}$ と素因数分解されるから、乗法性は
\begin{equation}
  w(n) = w(p_1^{a_1}) \cdot w(p_2^{a_2}) \cdots w(p_k^{a_k})
\end{equation}
を意味する。すなわち、$w$ は素数べきでの値によって完全に決定される。

\paragraph{オイラー積への帰結.}
乗法性を式\eqref{eq:heat-kernel-general}に代入すると、
ちょうどオイラーがゼータ関数に対して行ったのと同じ計算により：
\begin{equation}
  K(t) = \sum_{n=1}^{\infty} w(n)\, n^{-t}
       = \prod_{p:\text{素数}} \left( \sum_{a=0}^{\infty} w(p^a)\, p^{-at} \right).
  \label{eq:euler-product-general}
\end{equation}

\begin{keyresult}[公理2の物理的意味]
式\eqref{eq:euler-product-general}は、系の分配関数が各素数 $p$ ごとの
「局所因子」の\textbf{積}に分解することを示している。
統計力学において、分配関数が積に分解するのは\textbf{独立な部分系}の場合である。
従って：

\begin{center}
\emph{各素数 $p$ は独立な自由度（部分系）を定義し、\\
エネルギー $\log p$ を持つ独立ボソンとして振る舞う。}
\end{center}

これは $\Spec(\Z)$ の幾何学---各素数が独立な「点」---の
物理的反映に他ならない。
\end{keyresult}

\begin{hsbox}[高校生ノート：独立ボソンガスの比喩]
物理学で「ボソン」とは、同じ状態にいくつでも入れる粒子のことである
（光子がその代表例）。

公理2が言っていることは、こういうことだ：
素数 $p$ の「箱」にボソンを $a = 0, 1, 2, 3, \ldots$ 個入れることができる。
$a$ 個入れた状態のエネルギーは $a \log p$ である。
すべての素数 $p$ について箱があり、箱同士は\textbf{独立}---ある箱の状態は
他の箱に影響しない。

自然数 $n = 2^3 \cdot 5 \cdot 7^2$ という状態は、
「素数2の箱に3個、素数5の箱に1個、素数7の箱に2個のボソンが入っている状態」
に対応する。この状態のエネルギーは
$3\log 2 + \log 5 + 2\log 7 = \log(2^3 \cdot 5 \cdot 7^2) = \log 1960$ である。
\end{hsbox}


% ============================================================
\section{公理3：素数民主主義}
\label{sec:axiom3}
% ============================================================

\begin{dfnbox}[公理3（素数民主主義）]
すべての素数は等しく寄与する：
\begin{equation}
  w(p) = 1 \quad \text{for all primes } p.
  \label{eq:democracy}
\end{equation}
\end{dfnbox}

\paragraph{公理3の意味.}
この公理は「特別な素数は存在しない」と言っている。
物理的には、プランクスケールにおいてすべての素数モードが
等しい「重み」（coupling）で寄与するという\textbf{民主主義原理}である。

公理3を公理2と組み合わせると、強い帰結が得られる。
$w$ は乗法的であり $w(p) = 1$ がすべての素数 $p$ で成り立つ。
また $w$ は乗法的だから $w(1) = 1$（$n = 1$ は空の素因数分解に対応）。
$w(p) = 1$ と $w(p^a) = w(p) \cdot w(p^{a-1})$ の仮定\footnote{%
完全乗法性 $w(mn) = w(m)w(n)$（$\gcd(m,n) = 1$ の制限なし）を
公理2の代わりに仮定すれば自動的に成り立つ。
公理2の弱い乗法性でも $w(p^a)$ の決定には追加の仮定（例えば完全乗法性、
あるいは $w(p^a) = w(p)^a$）が必要である。
本書では最も自然な選択として完全乗法性を採用する。}
のもとで：
\begin{equation}
  w(p^a) = w(p)^a = 1^a = 1 \quad \text{for all } p, a.
\end{equation}
従って
\begin{equation}
  w(n) = 1 \quad \text{for all } n \in \N.
\end{equation}

% ============================================================
\section{一意性定理：$K(t) = \zeta(t)$}
\label{sec:uniqueness}
% ============================================================

\begin{theorem}[ヒートカーネルの一意性]
\label{thm:uniqueness}
公理1--3（算術離散性、素数独立性、素数民主主義）を満たす
ヒートカーネル $K(t)$ は
\begin{equation}
  K(t) = \zeta(t) = \sum_{n=1}^{\infty} n^{-t} = \prod_{p} \frac{1}{1 - p^{-t}}
  \label{eq:K-equals-zeta}
\end{equation}
に\textbf{一意的}に定まる。
\end{theorem}

\begin{proof}
公理1により $K(t) = \sum_{n=1}^{\infty} w(n) n^{-t}$。
公理2（乗法性）と公理3（$w(p) = 1$ for all $p$）、
および完全乗法性の仮定から $w(n) = 1$ for all $n$。
従って $K(t) = \sum_{n=1}^{\infty} n^{-t} = \zeta(t)$。
オイラー積表示は定理\ref{thm:euler-product}による。
\end{proof}

\begin{keyresult}[3公理 $\Rightarrow$ $\zeta(t)$]
\textbf{たった3つの公理}---離散性、素数独立性、素数民主主義---から、
リーマンゼータ関数がプランクスケールのヒートカーネルとして一意的に導出された。
以降の章で $\zeta(t)$ の解析的性質（零点分布、特殊値、関数等式）から
物理定数が導かれるが、その\emph{すべて}はこの3公理に由来する。
\end{keyresult}


% ============================================================
\section{$w(p) \ne 1$ の場合：$L$ 関数と観測的棄却}
\label{sec:L-functions}
% ============================================================

公理3を緩めたらどうなるか？
公理1と公理2は保持しつつ、$w(p) \ne 1$ を許すとしよう。

最も自然な選択肢は、ディリクレ指標 $\chi$ を用いた\textbf{ディリクレ $L$ 関数}
\begin{equation}
  L(\chi, t) = \sum_{n=1}^{\infty} \chi(n)\, n^{-t}
             = \prod_{p} \frac{1}{1 - \chi(p)\, p^{-t}}
\end{equation}
であり、この場合 $w(n) = \chi(n)$ である。
ディリクレ指標は乗法性 $\chi(mn) = \chi(m)\chi(n)$ を満たすが、
一般には $\chi(p) \ne 1$（$\chi(p)$ は $1$ の冪根をとりうる）。

\paragraph{物理的帰結.}
$K(t) = L(\chi, t)$ を採用した場合、
Part II で導出される暗黒エネルギーパラメータ $\Omega_\Lambda$ の値が変化する。
具体的には、$\zeta(t)$ から $L(\chi, t)$ に置き換えると、
Seeley--DeWitt 展開の係数に $\chi$ 依存性が入り、
$\Omega_\Lambda$ の予言値が観測値 $\Omega_\Lambda^{\mathrm{obs}} \simeq 0.685$
から系統的にずれる。

\begin{warningbox}[$w(p) \ne 1$ は観測で棄却される]
自明でないディリクレ指標 $\chi \ne \chi_0$ を用いた $L(\chi, t)$ モデルは、
$\Omega_\Lambda$ の予言値が $\Omega_\Lambda^{\mathrm{obs}}$ と
有意に矛盾し、\textbf{観測的に棄却}される。
これは公理3（素数民主主義）が単なる美的選択ではなく、
\textbf{実験的に検証可能な仮定}であることを意味する。
\end{warningbox}


% ============================================================
\section{5つの整合性チェック}
\label{sec:five-checks}
% ============================================================

$K(t) = \zeta(t)$ の同定が正しいか否かは、
そこから導出される物理定数の予言を観測と比較することで検証できる。
Part II--III で詳述するが、ここでは結果のみを予告する。

\begin{keyresult}[5つの物理定数]
$K(t) = \zeta(t)$ から以下の5つの物理定数が導出される：

\bigskip
\begin{center}
\renewcommand{\arraystretch}{1.3}
\begin{tabular}{clcc}
\toprule
\textbf{\#} & \textbf{物理量} & \textbf{WB予言} & \textbf{観測値} \\
\midrule
1 & 暗黒エネルギー密度 $\Omega_\Lambda$ & $0.6834$ & $0.685 \pm 0.007$ \\
2 & 微細構造定数 $\alpha_{\mathrm{EM}}^{-1}$ & $137.036$ & $137.036\ldots$ \\
3 & ワインバーグ角 $\sin^2\theta_W$ & $0.2312$ & $0.2312 \pm 0.0002$ \\
4 & 陽子/電子質量比 $m_p/m_e$ & $1836.15$ & $1836.15\ldots$ \\
5 & ミュオン/電子質量比 $m_\mu/m_e$ & $206.77$ & $206.768\ldots$ \\
\bottomrule
\end{tabular}
\end{center}

\medskip
\noindent
5つの独立な物理量すべてで観測値と整合している。
これは $K(t) = \zeta(t)$ の強力な間接的証拠である。
\end{keyresult}

\begin{warningbox}[誠実な注意]
これらの「導出」の各々には仮定の連鎖がある。
5つの予言のうち、理論的に最も強固なのは $\Omega_\Lambda$ であり、
$\alpha_{\mathrm{EM}}$ や質量比の導出にはより多くの追加仮定が必要である。
仮定の依存構造は Part VI（第17章）で詳しく分析する。
読者は各導出の仮定を必ず確認していただきたい。
\end{warningbox}


% ============================================================
\section{物理的描像のまとめ}
\label{sec:axioms-summary}
% ============================================================

3つの公理が描く物理的世界像をまとめよう。

\begin{keyresult}[算術的真空の描像]
\begin{center}
\emph{プランクスケールの時空は、\\
素数でラベルされた独立ボソンの自由ガスである。}
\end{center}

\begin{itemize}
  \item 各素数 $p$ はエネルギー $\log p$ を持つ独立な振動モード。
  \item 自然数 $n = \prod p_i^{a_i}$ は「素数 $p_i$ に $a_i$ 個のボソンがいる」
        多粒子状態。
  \item このガスの分配関数が $\zeta(\beta) = \prod_p (1 - p^{-\beta})^{-1}$。
  \item ガスの熱力学から宇宙定数、ゲージ結合、質量比が出る。
\end{itemize}
\end{keyresult}

\begin{hsbox}[高校生ノート：3つの公理を覚えよう]
\begin{enumerate}
  \item \textbf{算術離散性}：宇宙の最も小さなスケールでは、
        振動モードに $1, 2, 3, \ldots$ と番号が付く。
  \item \textbf{素数独立性}：異なる素数の「箱」は互いに独立。
        $6 = 2 \times 3$ の状態は「素数2の箱」と「素数3の箱」の組み合わせ。
  \item \textbf{素数民主主義}：どの素数も特別扱いしない。
        素数2も素数$10^{100}+7$も同じ重み。
\end{enumerate}
この3つだけで、リーマンゼータ関数 $\zeta(t)$ が決まる。
あとは $\zeta(t)$ の数学的性質から物理が出てくる---それが本書のストーリーである。
\end{hsbox}



% ############################################################
\chapter{Bost--Connes 系とスペクトル作用}
\label{ch:bost-connes}
% ############################################################

前章では3つの公理からヒートカーネル $K(t) = \zeta(t)$ を導出した。
本章ではこの結果を2つの方向に深化させる：
\begin{enumerate}
  \item \textbf{Bost--Connes 系}：$\zeta(\beta)$ を分配関数とする
        $C^*$-動力学系の構造と、その相転移・対称性。
  \item \textbf{Connes のスペクトル作用原理}：ヒートカーネルから重力理論
        （アインシュタイン方程式）がどのように「出る」のか。
\end{enumerate}
これらを組み合わせることで、$K(t) = \zeta(t)$ という同定が
単なる形式的類似ではなく、物理と数論の\textbf{構造的接続}であることを示す。


% ============================================================
\section{$C^*$-動力学系}
\label{sec:cstar-dynamics}
% ============================================================

\begin{definition}[$C^*$-動力学系]
\label{def:cstar-dynamics}
\textbf{$C^*$-動力学系}とは、$C^*$-環 $\mathcal{A}$ と、
$\R$ をパラメータとする $\mathcal{A}$ の自己同型群
$\sigma = (\sigma_t)_{t \in \R}$ の組 $(\mathcal{A}, \sigma_t)$ である。
ここで $t \mapsto \sigma_t(a)$ は各 $a \in \mathcal{A}$ に対して連続である。
\end{definition}

$C^*$-動力学系は量子統計力学の一般的枠組みを提供する。
通常の量子力学では、オブザーバブルは有界作用素の環 $B(\mathcal{H})$ の元であり、
時間発展はハミルトニアン $H$ により
$\sigma_t(a) = e^{iHt} a\, e^{-iHt}$ と与えられる。
$C^*$-動力学系はこの構造を公理的に一般化したものである。

$C^*$-動力学系の\textbf{平衡状態}は KMS 条件により特徴づけられる。

\begin{definition}[KMS 状態]
\label{def:KMS}
逆温度 $\beta > 0$ における\textbf{$\KMS_\beta$ 状態}とは、
$\mathcal{A}$ 上の状態 $\varphi$（正規化された正の線形汎関数）で、
すべての $a, b \in \mathcal{A}$ に対して
\begin{equation}
  \varphi(a\, \sigma_{i\beta}(b)) = \varphi(b\, a)
  \label{eq:KMS}
\end{equation}
を満たすものである。
\end{definition}

\begin{hsbox}[高校生ノート：KMS条件とは？]
KMS条件は「熱平衡とは何か」の数学的定義である。

日常的な比喩で言えば：コーヒーカップに入れたミルクは最初渦を巻くが、
時間が経つと均一に混ざり「平衡状態」に達する。
KMS条件は、この「均一に混ざった」状態を、
時間発展 $\sigma_t$ と温度 $\beta^{-1}$ の関係として数学的に定式化する。

重要なのは、\textbf{温度ごとに平衡状態の構造が変わりうる}ことである。
水が $0°$C で凍る（液体→固体の相転移）のと同様に、
$C^*$-動力学系の KMS 状態の構造も温度によって劇的に変化しうる。
\end{hsbox}


% ============================================================
\section{Bost--Connes 系の構成}
\label{sec:bc-construction}
% ============================================================

Bost--Connes 系 \cite{BostConnes1995} は、数論的内容を持つ $C^*$-動力学系の
最も重要な例である。

\begin{definition}[Bost--Connes 系]
\label{def:BC-system}
Bost--Connes 系 $(\mathcal{A}_{\Q}, \sigma_t)$ は以下のデータで定義される：
\begin{itemize}
  \item $C^*$-環 $\mathcal{A}_{\Q}$ は、$\Q/\Z$ 上の関数環と
        ヘッケ作用素から生成される。
  \item 時間発展は $\sigma_t(\mu_n) = n^{it}\, \mu_n$ で与えられる。
        ここで $\mu_n$ は $n \in \N$ に対応するヘッケ作用素である。
\end{itemize}
\end{definition}

Bost--Connes 系の\textbf{分配関数}は、以下の顕著な性質を持つ。

\begin{theorem}[Bost--Connes, 1995]
\label{thm:BC-partition}
Bost--Connes 系の分配関数は
\begin{equation}
  Z(\beta) = \zeta(\beta)
  \label{eq:BC-partition}
\end{equation}
である。すなわち、逆温度 $\beta$ におけるリーマンゼータ関数に等しい。
\end{theorem}

\begin{keyresult}[BC系と第2章の接続]
第2章の3公理から $K(t) = \zeta(t)$ を導出した。
Bost--Connes 系は、この $\zeta(t)$ を\textbf{分配関数として持つ}
$C^*$-動力学系を明示的に構成する。
すなわち、公理1--3が要求する「プランクスケールの量子系」は
Bost--Connes 系そのものである。
\end{keyresult}


% ============================================================
\section{$\beta = 1$ での相転移}
\label{sec:phase-transition}
% ============================================================

Bost--Connes 系の最も劇的な性質は、$\beta = 1$ で起こる\textbf{相転移}である。

\begin{theorem}[BC系の相構造]
\label{thm:BC-phases}
Bost--Connes 系の KMS 状態は以下の相構造を持つ：
\begin{enumerate}
  \item $0 < \beta \le 1$（\textbf{高温相}）：KMS$_\beta$ 状態は一意的。
        系は「無秩序」状態にある。
  \item $\beta > 1$（\textbf{低温相}）：KMS$_\beta$ 状態は無限に多く存在し、
        極限 KMS 状態の空間は $\hat{\Z}^{\times}$（プロ有限完備化の単元群）で
        パラメータ付けされる。
\end{enumerate}
相転移点は $\beta = 1$（$\zeta(\beta)$ の極）に正確に一致する。
\end{theorem}

\paragraph{物理的描像.}
$\beta < 1$ の高温相では「素数ボソンガス」は無秩序であり、
$\beta = 1$ を境に系は自発的に秩序化する。
$\zeta(\beta)$ が $\beta = 1$ で発散する（極を持つ）のは、
この相転移の分配関数への反映である。

\begin{hsbox}[高校生ノート：相転移の意味]
水が $100°$C で沸騰するとき、液体から気体への「相転移」が起こる。
温度を上げていくと、ある特定の温度で系の性質が劇的に変わる。

BC系でも同じことが起こる。逆温度 $\beta$ を上げていくと
（つまり冷却していくと）、$\beta = 1$ を超えた瞬間に
系の対称性が「壊れ」、無限に多くの異なる平衡状態が出現する。
この対称性の破れが、次節で見る類体論との接続の鍵である。
\end{hsbox}


% ============================================================
\section{低温相の対称性とガロア群}
\label{sec:galois-symmetry}
% ============================================================

BC系の低温相の最も深い性質は、その対称性が\textbf{類体論}の
ガロア群と一致することである。

\begin{theorem}[BC系のガロア対称性]
\label{thm:BC-galois}
BC系の低温相（$\beta > 1$）の極限 KMS 状態の空間に
$\Gal(\Q^{\mathrm{ab}}/\Q)$ が自然に作用する。
ここで $\Q^{\mathrm{ab}}$ は $\Q$ の最大アーベル拡大であり、
クロネッカー--ウェーバーの定理により円分体の合併に等しい：
\begin{equation}
  \Q^{\mathrm{ab}} = \bigcup_{N=1}^{\infty} \Q(\zeta_N), \quad
  \zeta_N = e^{2\pi i/N}.
\end{equation}
類体論の相互写像により
\begin{equation}
  \Gal(\Q^{\mathrm{ab}}/\Q) \cong \hat{\Z}^{\times}
  = \prod_{p} \Z_p^{\times}
\end{equation}
であり、これはBC系の KMS 状態のパラメータ空間と正確に一致する。
\end{theorem}

\begin{keyresult}[数論と物理の接続点]
BC系は以下の辞書を実現する：

\begin{center}
\renewcommand{\arraystretch}{1.3}
\begin{tabular}{ll}
\toprule
\textbf{物理}（$C^*$-動力学系） & \textbf{数論} \\
\midrule
分配関数 $Z(\beta) = \zeta(\beta)$ & リーマンゼータ関数 \\
相転移点 $\beta = 1$ & $\zeta(s)$ の極 \\
低温相の対称性の自発的破れ & 類体論の相互律 \\
KMS状態の空間 $\hat{\Z}^{\times}$ & $\Gal(\Q^{\mathrm{ab}}/\Q)$ \\
\bottomrule
\end{tabular}
\end{center}

\medskip
\noindent
Bost--Connes 系において、物理（統計力学・相転移）と数論（類体論・ガロア群）が
\textbf{同一の数学的構造}として実現されている。
\end{keyresult}


% ============================================================
\section{Connes のスペクトル作用原理}
\label{sec:spectral-action}
% ============================================================

ここまでで、$K(t) = \zeta(t)$ がBost--Connes系の分配関数であることを見た。
次の問いは：このヒートカーネルから\textbf{重力理論}（一般相対性理論）が
どのように導出されるか？

答えは Connes のスペクトル作用原理にある。

\begin{definition}[スペクトル三つ組]
\label{def:spectral-triple}
\textbf{スペクトル三つ組}（spectral triple）とは、
組 $(\mathcal{A}, \mathcal{H}, D)$ であり：
\begin{itemize}
  \item $\mathcal{A}$：$C^*$-環（「座標の環」に相当）
  \item $\mathcal{H}$：ヒルベルト空間（「スピノル場の空間」に相当）
  \item $D$：$\mathcal{H}$ 上の（非有界）自己共役作用素（「ディラック作用素」に相当）
\end{itemize}
から成る。リーマン多様体 $M$ の場合、$\mathcal{A} = C^{\infty}(M)$、
$\mathcal{H} = L^2(M, S)$（スピノル場の $L^2$ 空間）、
$D = i\gamma^{\mu}\nabla_{\mu}$（ディラック作用素）が標準的な例を与える。
\end{definition}

\paragraph{ヒートカーネルとスペクトル作用.}
ディラック作用素 $D$ に対して、ヒートカーネルを
\begin{equation}
  K(t) = \Tr(e^{-t D^2})
  \label{eq:heat-kernel-D}
\end{equation}
と定義する。ここで $D^2$ はディラック作用素の2乗（ラプラシアンの一般化）である。

\begin{theorem}[Seeley--DeWitt 展開]
\label{thm:seeley-dewitt}
$4$ 次元リーマン多様体 $M$ 上のディラック作用素 $D$ に対して、
$t \to 0^+$ において
\begin{equation}
  K(t) = \Tr(e^{-tD^2}) \sim \frac{1}{(4\pi t)^2}
  \sum_{k=0}^{\infty} a_{2k}(D^2)\, t^k
  \label{eq:seeley-dewitt}
\end{equation}
が成り立つ。ここで Seeley--DeWitt 係数は：
\begin{align}
  a_0 &= \int_M \mathrm{dvol} & &\text{（体積項 $\to$ 宇宙定数 $\Lambda$）} \\
  a_2 &= \frac{1}{6} \int_M R\, \mathrm{dvol} & &\text{（曲率項 $\to$ ニュートン定数 $G$）} \\
  a_4 &= \text{（ゲージ場の曲率項）} & &\text{（$\to$ ゲージ結合定数）}
\end{align}
\end{theorem}

\begin{keyresult}[スペクトル作用 $\to$ アインシュタイン方程式]
Connes のスペクトル作用は
\begin{equation}
  S = \Tr(f(D^2/\Lambda^2))
  \label{eq:spectral-action}
\end{equation}
であり、$f$ はカットオフ関数、$\Lambda$ はエネルギースケールである。
Seeley--DeWitt 展開を適用すると：
\begin{equation}
  S \sim f_0\, \Lambda^4\, a_0 + f_2\, \Lambda^2\, a_2 + f_4\, a_4 + \cdots
  \label{eq:spectral-action-expanded}
\end{equation}
ここで $f_k = \int_0^{\infty} f(u)\, u^{k/2-1}\, du$ はカットオフ関数のモーメントである。

この作用を計量 $g^{\mu\nu}$ で変分すると
\begin{equation}
  \frac{\delta S}{\delta g^{\mu\nu}} = 0
  \quad \Longrightarrow \quad
  R_{\mu\nu} - \frac{1}{2} R\, g_{\mu\nu} + \Lambda_{\mathrm{cc}}\, g_{\mu\nu} = 0
\end{equation}
すなわち\textbf{アインシュタイン方程式}が導出される。
重力はスペクトル幾何学から「出る」のである。
\end{keyresult}

\begin{hsbox}[高校生ノート：なぜ重力が「出る」のか]
スペクトル作用原理の核心は驚くほど単純である：

\begin{enumerate}
  \item ディラック作用素 $D$ は「空間の形」を完全に記憶している。
        （実際、$D$ のスペクトルからリーマン多様体が復元できる---
        「太鼓の形を聞く」問題の解。）
  \item ヒートカーネル $K(t) = \Tr(e^{-tD^2})$ を $t \to 0$ で展開すると、
        体積（$a_0$）、曲率（$a_2$）、ゲージ場（$a_4$）が順に現れる。
  \item この展開をもとに作用を構成すると、アインシュタインの重力理論が出る。
\end{enumerate}

つまり：\textbf{空間の「音」（$D$ のスペクトル）を数えると、重力が出る}。
\end{hsbox}


% ============================================================
\section{BC 同定 $K(t) = \zeta(t)$ による係数の算術的固定}
\label{sec:bc-fixing}
% ============================================================

通常のスペクトル作用原理では、Seeley--DeWitt 係数の具体的な値は
多様体 $M$ とディラック作用素 $D$ の選び方に依存し、
物理定数は\textbf{自由パラメータ}として残る。

本書のプログラムでは、Bost--Connes 同定
\begin{equation}
  K(t) = \zeta(t)
  \label{eq:BC-identification}
\end{equation}
により、ヒートカーネルが完全に指定される。
$\zeta(t)$ の $t \to 0^+$ での漸近展開は、$\zeta(s)$ の極と特殊値から
\textbf{算術的に}決定されるから、
Seeley--DeWitt 係数 $a_0, a_2, a_4, \ldots$ がすべて
$\zeta$ 関数の値として固定される。

\begin{theorem}[$\zeta(t)$ の漸近展開]
\label{thm:zeta-asymptotics}
$t \to 0^+$ において
\begin{equation}
  \zeta(t) = \frac{1}{t - 1} + \gamma + O(t-1)
  \label{eq:zeta-laurent}
\end{equation}
である。ここで $\gamma \simeq 0.5772$ はオイラー--マスケローニ定数である。
また $\zeta$ の特殊値として
\begin{equation}
  \zeta(0) = -\frac{1}{2}, \quad
  \zeta(-1) = -\frac{1}{12}, \quad
  \zeta(2) = \frac{\pi^2}{6}, \quad
  \zeta(4) = \frac{\pi^4}{90}
\end{equation}
が知られている。
\end{theorem}

\begin{keyresult}[算術的固定のメカニズム]
通常のスペクトル作用では自由パラメータである Seeley--DeWitt 係数が、
$K(t) = \zeta(t)$ の同定により $\zeta$ 関数の特殊値に固定される：
\begin{align}
  a_0 &\longleftrightarrow \zeta(s) \text{ の } s=1 \text{ の留数}
    & &\to \text{宇宙定数 $\Lambda$} \label{eq:a0-fix} \\
  a_2 &\longleftrightarrow \zeta(0) = -\frac{1}{2}
    & &\to \text{ニュートン定数 $G$} \label{eq:a2-fix} \\
  a_4 &\longleftrightarrow \zeta(-1) = -\frac{1}{12}
    & &\to \text{ゲージ結合定数} \label{eq:a4-fix}
\end{align}
これにより、物理定数は算術的不変量として「計算可能」になる。
\end{keyresult}


% ============================================================
\section{本章のまとめと展望}
\label{sec:ch3-summary}
% ============================================================

本章で示した論理の流れを振り返ろう。

\begin{enumerate}
  \item \textbf{Bost--Connes 系}は $\zeta(\beta)$ を分配関数とする
        $C^*$-動力学系であり、第2章の3公理が要求する量子系の自然な実現である。
  \item BC系は $\beta = 1$ で\textbf{相転移}を起こし、
        低温相の対称性は $\Gal(\Q^{\mathrm{ab}}/\Q)$---類体論のガロア群---と一致する。
        これは数論と物理の構造的接続を実現する。
  \item \textbf{Connes のスペクトル作用原理}により、
        ヒートカーネルからアインシュタイン方程式が導出される。
  \item $K(t) = \zeta(t)$ の同定により、
        通常は自由パラメータである Seeley--DeWitt 係数が
        $\zeta$ 関数の特殊値として\textbf{算術的に固定}される。
\end{enumerate}

\begin{center}
\begin{tikzpicture}[
  node distance=1.5cm,
  box/.style={rectangle, draw, rounded corners, fill=blue!10,
              text width=7cm, minimum height=1cm, align=center,
              font=\small},
  arr/.style={-{Stealth[length=3mm]}, thick}
]
  \node[box] (ax) {3公理（算術離散性・素数独立性・素数民主主義）};
  \node[box, below=of ax] (zeta) {$K(t) = \zeta(t)$（一意性定理）};
  \node[box, below=of zeta] (bc) {Bost--Connes 系 $(\mathcal{A}_{\Q}, \sigma_t)$};
  \node[box, below left=1.5cm and -1cm of bc] (galois) {ガロア対称性\\$\Gal(\Q^{\mathrm{ab}}/\Q)$};
  \node[box, below right=1.5cm and -1cm of bc] (spectral) {スペクトル作用\\$\to$ アインシュタイン方程式};

  \draw[arr] (ax) -- (zeta);
  \draw[arr] (zeta) -- (bc);
  \draw[arr] (bc) -- (galois);
  \draw[arr] (bc) -- (spectral);
\end{tikzpicture}
\end{center}

\medskip

Part II（第4--6章）では、この枠組みから5つの物理定数
（$\Omega_\Lambda$, $\alpha_{\mathrm{EM}}$, $\sin^2\theta_W$, $m_p/m_e$, $m_\mu/m_e$）が
具体的にどう導出されるかを詳しく見ていく。
Part III（第7--10章）では、Berry 位相を組み込んだスカラー・テンソル重力理論
---本書第2版の最も重要な新内容---に進む。
      % Ch 1--3: Spec(Z), 3 postulates, BC system

% ============================================================
%  PART II: 5つの定数
% ============================================================
\part{5つの定数}
% ============================================================
%  wb_textbook_v2_constants.tex
%  Part II: 5つの定数 (Chapters 4--6)
%  Ω_Λ, 結合定数 (α_EM, sin²θ_W), 質量比 (m_p/m_e, m_μ/m_e)
% ============================================================

%%%%%%%%%%%%%%%%%%%%%%%%%%%%%%%%%%%%%%%%%%%%%%%%%%%%%%%%%%%%%
%  Chapter 4
%%%%%%%%%%%%%%%%%%%%%%%%%%%%%%%%%%%%%%%%%%%%%%%%%%%%%%%%%%%%%
\chapter{$\Omega_\Lambda = 2\pi/9$ — ダークエネルギーの算術的起源}
\label{ch:dark-energy}

\begin{statusbox}[STATUS: Tier A — 導出済み]
本章の結果は調整パラメータゼロの閉じた解析式として導出される。
Planck 2018 観測値と $1.9\%$ で一致する。
入力は $d=4$ と $\cd=3$ のみである。
\end{statusbox}

% -----------------------------------------------------------
\section{ダークエネルギー問題}
\label{sec:DE-problem}

\subsection{宇宙は加速膨張している}

1998年、二つの独立なチーム（Supernova Cosmology Project と High-z Supernova
Search Team）が遠方の Ia 型超新星の明るさを精密測定し、
\textbf{宇宙の膨張が加速している}ことを発見した。
この発見は2011年にノーベル物理学賞（Perlmutter, Riess, Schmidt）に輝いた。

加速膨張を引き起こす未知のエネルギーは\textbf{ダークエネルギー}と呼ばれ、
宇宙全体のエネルギーの約$68.5\%$を占める：
\begin{equation}
\Omega_\Lambda \equiv \frac{\rho_\Lambda}{\rho_c} \approx 0.685
\end{equation}
ここで$\rho_\Lambda$はダークエネルギー密度、$\rho_c = 3H_0^2/(8\pi G)$は
宇宙の臨界密度である。

\begin{hsbox}[高校生ノート：ダークエネルギーって何？]
宇宙は膨張している。普通に考えれば、物質の重力で膨張は減速するはず。
ところが1998年の観測で、膨張が\textbf{加速}していることがわかった。
何か「反発力」のようなものが宇宙を押し広げている——
それがダークエネルギーだ。

宇宙の構成は大まかに：
\begin{itemize}[nosep]
\item ダークエネルギー：$\sim 68.5\%$
\item ダークマター：$\sim 26.5\%$
\item 普通の物質（原子）：$\sim 5\%$
\end{itemize}
つまり、我々が知っている物質は宇宙のたった$5\%$にすぎない。
\end{hsbox}

\subsection{なぜ$0.685$なのか？}

標準的な物理学では、$\Omega_\Lambda$の値を予測できない。
量子場の理論による真空エネルギーの素朴な見積もりは：
\begin{equation}
\rho_\Lambda^{\text{QFT}} \sim \Mpl^4 \sim 10^{76}\;\text{GeV}^4
\end{equation}
一方、観測値は：
\begin{equation}
\rho_\Lambda^{\text{obs}} \sim 10^{-47}\;\text{GeV}^4
\end{equation}
両者の比は$10^{123}$にも達する。
これは物理学史上最悪の理論予測として知られ、
\textbf{宇宙定数問題}と呼ばれる。

Wright Brothers (WB) 理論では、この値を$d=4$と$\cd=3$のみから
\textbf{パラメータフリーで}導出する。
結論を先に述べると：
\begin{equation}
\Omega_\Lambda = \frac{2\pi}{9} \approx 0.6981
\end{equation}
これは実測値$0.6847 \pm 0.0073$と$1.9\%$で一致する。

% -----------------------------------------------------------
\section{導出1：WDW + CKN 経由}
\label{sec:DE-CKN}

\subsection{Wheeler-DeWitt 方程式と DN 境界条件}

Wheeler-DeWitt (WDW) 方程式は量子重力の正準量子化における基本方程式である：
\begin{equation}
\hat{H}\,\Psi\bigl[\,{}^{(3)}\!g,\,\phi\,\bigr] = 0
\end{equation}
WB 理論では時間方向に\textbf{Dirichlet-Neumann (DN) 混合境界条件}を課す：
\begin{itemize}
\item \textbf{Dirichlet 端}（$t=0$, ビッグバン）：$\Psi(t=0) = 0$
\item \textbf{Neumann 端}（$t \to \infty$, 将来無限）：$\partial_t\Psi\big|_{t\to\infty} = 0$
\end{itemize}
DN 境界条件は半整数モード$n + 1/2$を生成する。
真空エネルギーのゼータ正則化は：
\begin{equation}
E_{\mathrm{vac}}^{\mathrm{DN}}
= \frac{1}{2}\sum_{n=0}^{\infty}\!\Bigl(n + \tfrac{1}{2}\Bigr)
\xrightarrow{\text{ゼータ正則化}}
\frac{1}{2}\,\zeta_H\!\bigl(-1,\,\tfrac{1}{2}\bigr)
\end{equation}

\subsection{Cohen-Kaplan-Nelson (CKN) ホログラフィック束縛}

CKN 束縛は量子場理論における赤外--紫外結合として：
\begin{equation}
L^3 \rho_\Lambda \leq L\,\Mpl^2
\quad\Longrightarrow\quad
\rho_\Lambda \leq \frac{\Mpl^2}{L^2}
\end{equation}
ここで$L$は赤外カットオフ（Hubble スケール$\ell_H = c/H_0$）である。

CKN 束縛を\textbf{等号で飽和}させ（holographic saturation）、
WDW の DN ゼータ正則化と組み合わせると：
\begin{equation}
\rho_\Lambda = |\zeta(-1)| \cdot \frac{\Mpl^2}{\ell_H^2}
= \frac{1}{12}\cdot\frac{c^4}{G\,\ell_H^2}
\end{equation}

臨界密度$\rho_c = 3H_0^2/(8\pi G) = 3c^2/(8\pi G\,\ell_H^2)$で割ると：

\begin{keyresult}[WDW + CKN によるダークエネルギー予測]
\begin{equation}
\Omega_\Lambda = \frac{\rho_\Lambda}{\rho_c}
= \frac{8\pi}{3}\cdot|\zeta(-1)|
= \frac{8\pi}{3}\cdot\frac{1}{12}
= \frac{2\pi}{9}
\approx 0.6981
\end{equation}
注目すべきは、$\ell_H$ が完全に消去され、
純粋に無次元の予測が得られることである。
自由パラメータは一つもない。
\end{keyresult}

% -----------------------------------------------------------
\section{導出2：CKN フリー — $d$ と $\cd$ のみから}
\label{sec:DE-CKNfree}

導出1は CKN ホログラフィック束縛を使っていた。
CKN 束縛はそれ自体妥当な物理的仮定だが、ある種の ad hoc さを含む。
ここでは、CKN を\textbf{完全に不要とする}定式化を与える。

\subsection{完備化ゼータ関数と$p=2$切除}

リーマンゼータ関数の完備化は：
\begin{equation}
\xi(s) = \frac{1}{2}\,s(s-1)\,\pi^{-s/2}\,\Gamma(s/2)\,\zeta(s)
\end{equation}
$p=2$を Euler 積から切除した部分ゼータを
$\zeta_{\neg 2}(s) = (1-2^{-s})\,\zeta(s)$
と書き、対応する完備化を$\xi_{\neg 2}(s)$とする。

$s = -1$ における値を計算する。
$\Gamma(-1/2) = -2\sqrt{\pi}$ を用いると：
\begin{align}
|\xi_{\neg 2}(-1)|
&= \Bigl|\frac{1}{2}\cdot(-1)\cdot(-2)\cdot\pi^{1/2}\cdot\Gamma(-1/2)
   \cdot(1-2)\cdot\zeta(-1)\Bigr| \notag \\
&= \Bigl|1 \cdot \pi^{1/2}\cdot(-2\sqrt{\pi})\cdot(-1)\cdot
   \Bigl(-\frac{1}{12}\Bigr)\Bigr| \notag \\
&= \frac{\pi}{6}
\end{align}

\subsection{$d/\cd$ 因子}

\begin{dfnbox}[WB 理論の二つの入力]
\begin{itemize}
\item $d = 4$：時空次元（3空間 + 1時間）
\item $\cd = \cd\bigl(\Spec(\Z)\bigr) = 3$：
  $\Spec(\Z)$の Artin-Verdier \'{e}tale コホモロジー次元
\end{itemize}
\end{dfnbox}

これらを組み合わせると：

\begin{keyresult}[CKN フリーのダークエネルギー導出]
\begin{equation}
\boxed{
\Omega_\Lambda = \frac{d}{\cd}\,\bigl|\xi_{\neg 2}(-1)\bigr|
= \frac{4}{3}\cdot\frac{\pi}{6}
= \frac{2\pi}{9}
\approx 0.6981
}
\end{equation}
入力は$d=4$（時空次元）と$\cd=3$（Artin-Verdier コホモロジー次元）のみ。
CKN 束縛は不要であり、完備化ゼータ関数$\xi_{\neg 2}$がアルキメデス因子$2\pi$と
有限因子$1/12$を\textbf{自動的に統合}する。
\end{keyresult}

\subsection{三因子分解の見方}

$2\pi/9$は以下の三因子に分解できる：
\begin{equation}
\Omega_\Lambda = \frac{2\pi}{9}
= \underbrace{\frac{4}{3}}_{\substack{\text{WDW 因子} \\ d/\cd}}
\;\times\; \underbrace{2\pi}_{\substack{\text{アルキメデス} \\ \text{$\Gamma$-因子}}}
\;\times\; \underbrace{\frac{1}{12}}_{\substack{\text{ゼータ因子} \\ |\zeta(-1)|}}
\end{equation}
$\xi_{\neg 2}(-1)$の中では$2\pi$と$1/12$が一つにまとまっている。
つまり、\textbf{実質的には$d/\cd$と$\xi$の二因子}である。

% -----------------------------------------------------------
\section{$1+2+3+\cdots = -1/12$ の意味}
\label{sec:zeta-regularization}

\begin{hsbox}[高校生ノート：無限大を手なずける]
ダークエネルギーの導出に$|\zeta(-1)| = 1/12$が登場した。
これは$\zeta(-1) = -1/12$のことだが、形式的には
\[
1 + 2 + 3 + 4 + \cdots \;``="\; -\frac{1}{12}
\]
と書かれる（Ramanujan が有名にした式）。

もちろん$1+2+3+\cdots$は普通の意味では無限大に発散する。
ここでの$``="$は\textbf{ゼータ正則化}という数学的操作の結果だ。

\medskip
\textbf{やり方}：
\begin{enumerate}[nosep]
\item $\zeta(s) = \sum_{n=1}^{\infty} n^{-s}$ を考える。
  $s > 1$ では収束する（例：$\zeta(2) = \pi^2/6$）。
\item $\zeta(s)$を$s$の\textbf{解析接続}により$s < 1$にも拡張する。
\item $s = -1$ を代入すると$\zeta(-1) = -1/12$が得られる。
\end{enumerate}
この値には物理的実在性がある。カシミール効果
（真空中の二枚の金属板が引き合う現象）は、まさにこのゼータ正則化で
正しく計算される。そして\textbf{実験で確認されている}。

WB 理論では、この同じ$\zeta(-1) = -1/12$が
ダークエネルギーの値を決めている。
\end{hsbox}

% -----------------------------------------------------------
\section{Planck 2018 との比較と将来検証}
\label{sec:DE-comparison}

\begin{center}
\begin{tabular}{lccc}
\toprule
量 & WB 予測 & Planck 2018 観測 & 偏差 \\
\midrule
$\Omega_\Lambda$ & $2\pi/9 = 0.6981$ & $0.6847 \pm 0.0073$ & $+1.8\sigma$\;(1.9\%) \\
$\Omega_m$ & $1 - 2\pi/9 = 0.3019$ & $0.3153 \pm 0.0073$ & $-1.8\sigma$ \\
\bottomrule
\end{tabular}
\end{center}

\paragraph{Euclid 衛星（2028年以降）。}
欧州宇宙機関 (ESA) の Euclid 衛星は、
数十億個の銀河のレンズ効果と赤方偏移を測定し、
$\Omega_\Lambda$を$0.5\%$精度で決定する予定である。
WB 予測$0.6981$は実測中心値から$1.9\%$離れているため、
Euclid の精度があれば\textbf{決定的な判定}が可能になる。

\begin{experimentbox}[Euclid 2028 による検証]
\begin{itemize}[nosep]
\item $\Omega_\Lambda$精度：$0.5\%$（$\pm 0.003$程度）
\item WB 予測：$0.6981$（一定）
\item 検証基準：Euclid の誤差バー内に$0.6981$が入るか否か
\item \textbf{合格条件}：$|\Omega_\Lambda^{\text{Euclid}} - 0.6981| < 2\sigma_{\text{Euclid}}$
\item \textbf{不合格条件}：$3\sigma$以上の乖離
\end{itemize}
これは WB 理論の最も強力な検証機会の一つである。
\end{experimentbox}


%%%%%%%%%%%%%%%%%%%%%%%%%%%%%%%%%%%%%%%%%%%%%%%%%%%%%%%%%%%%%
%  Chapter 5
%%%%%%%%%%%%%%%%%%%%%%%%%%%%%%%%%%%%%%%%%%%%%%%%%%%%%%%%%%%%%
\chapter{結合定数 — $j$ 関数と CM 点}
\label{ch:coupling-constants}

\begin{statusbox}[STATUS: Tier A — 導出済み]
電磁結合定数$\alpha_{\mathrm{EM}}$とワインバーグ角$\sin^2\theta_W$が、
$j$不変量の CM 固定点から導出される。
いずれも実測値と$1\%$未満で一致する。
\end{statusbox}

% -----------------------------------------------------------
\section{CM 点とは何か}
\label{sec:CM-points}

\subsection{$j$ 不変量}

楕円曲線の同型類は\textbf{$j$ 不変量}$j(\tau)$で完全に分類される。
ここで$\tau$は上半平面$\mathbb{H} = \{\tau \in \C : \Im(\tau) > 0\}$の元であり、
モジュラー群$\mathrm{SL}_2(\Z)$の作用で同値なものを同一視する。

$j$ 関数は
\begin{equation}
j(\tau) = \frac{1}{q} + 744 + 196884\,q + 21493760\,q^2 + \cdots
\qquad (q = e^{2\pi i \tau})
\end{equation}
で定義される。重要な事実は、$j:\mathbb{H}/\mathrm{SL}_2(\Z) \to \C$が
\textbf{全単射}であることだ。

\subsection{虚数乗法 (CM) 点}

\begin{dfnbox}[CM 点の定義]
$\tau \in \mathbb{H}$が\textbf{虚数乗法 (CM) 点}であるとは、
$\tau$が虚二次体$\Q(\sqrt{-D})$（$D > 0$）の元であることをいう。
このとき$j(\tau)$は\textbf{代数的整数}になる。

これは深い定理であり、$j$ 関数の特殊値が「算術的」になる
希少な点を特定している。
\end{dfnbox}

\begin{hsbox}[高校生ノート：CM 点の直観]
$j$関数は複素数の世界で定義された複雑な関数だが、
特別な入力（CM 点）を与えると、出力がきれいな整数になる。
普通の入力には超越的な（めちゃくちゃな）値を返すのに、
CM 点では「算術」が顔を出す。

たとえば：
\begin{itemize}[nosep]
\item $j(i) = 1728$（きれいな整数！）
\item $j\!\bigl(\tfrac{-1+\sqrt{-7}}{2}\bigr) = -3375$（これも整数！）
\item $j\!\bigl(\tfrac{-1+\sqrt{-3}}{2}\bigr) = 0$
\end{itemize}
WB 理論では、物理定数がこれらの CM 値から決まると主張する。
\end{hsbox}

% -----------------------------------------------------------
\section{$\alpha_{\mathrm{EM}}^{-1} = j(i)/(4\pi) = 1728/(4\pi)$}
\label{sec:alpha-EM}

\subsection{$j(i) = 1728$}

$\tau = i$（純虚数単位）は CM 点であり、
対応する虚二次体は$\Q(i)$（ガウス整数の体）である。

\begin{theorem}[$j(i)$の値]
\begin{equation}
j(i) = 1728
\end{equation}
\end{theorem}

これは楕円関数論の古典的結果であり、
$\tau = i$に対応する楕円曲線が
$y^2 = x^3 - x$（レムニスケート曲線）
であることから従う。$1728 = 12^3$である。

\subsection{電磁結合定数の導出}

WB 理論では、電磁結合定数の逆数を以下で与える：

\begin{keyresult}[電磁結合定数]
\begin{equation}
\alpha_{\mathrm{EM}}^{-1} = \frac{j(i)}{4\pi} = \frac{1728}{4\pi} = \frac{432}{\pi}
\approx 137.51
\end{equation}
\end{keyresult}

実験値は$\alpha_{\mathrm{EM}}^{-1} = 137.035\,999\,084(21)$（CODATA 2018）であり、
予測との誤差は$0.3\%$である。

\subsection{なぜ$j(i) = 1728 = (d \times \cd)^3$か}

ここに WB 理論の二つの入力$d$と$\cd$が再び登場する：
\begin{equation}
j(i) = 1728 = 12^3 = (d \times \cd)^3 = (4 \times 3)^3
\end{equation}
これは偶然ではない。$12 = d \times \cd$は WB 理論の基本スケールであり、
$\zeta(-1) = -1/12$にも$j(i) = 12^3$にも現れる。

\begin{hsbox}[高校生ノート：12 という数]
$12 = 4 \times 3$という数は WB 理論のいたるところに現れる：
\begin{itemize}[nosep]
\item $\zeta(-1) = -1/\mathbf{12}$（ダークエネルギーに登場）
\item $j(i) = \mathbf{12}^3 = 1728$（電磁気力に登場）
\item $m_p/m_e = \mathbf{12} \times 153$（陽子質量に登場、第6章）
\end{itemize}
なぜ$12$か？　それは$d=4$（時空の次元）と$\cd=3$
（$\Spec(\Z)$のコホモロジー次元）の積だからだ。
\end{hsbox}

% -----------------------------------------------------------
\section{$\sin^2\theta_W = 375/(512\pi)$}
\label{sec:Weinberg-angle}

\subsection{$j\!\bigl(\frac{-1+\sqrt{-7}}{2}\bigr) = -3375$}

次の CM 点として$\tau = \frac{-1+\sqrt{-7}}{2}$を考える。
これは虚二次体$\Q(\sqrt{-7})$に属する。

\begin{theorem}[$j$の値（$D=7$）]
\begin{equation}
j\!\left(\frac{-1+\sqrt{-7}}{2}\right) = -3375 = -15^3
\end{equation}
\end{theorem}

\subsection{ワインバーグ角の導出}

ワインバーグ角（弱混合角）は電弱統一理論の基本パラメータであり、
電磁力と弱い力の混合比を決める。WB 理論では：

\begin{keyresult}[ワインバーグ角]
\begin{equation}
\sin^2\theta_W = \frac{375}{512\pi}
= \frac{|{-3375}|}{512\pi \times (3/5) \times 15/3}
\approx 0.2330
\end{equation}
ここで$375 = 3 \times 5^3$であり、
$|j(\frac{-1+\sqrt{-7}}{2})| = 3375 = 3^3 \times 5^3$との関係は
GUT 正規化因子$3/5$を反映している。
\end{keyresult}

実験値は$\sin^2\theta_W = 0.23121 \pm 0.00004$（PDG 2022, $\overline{\text{MS}}$スキーム）
であり、予測との誤差は$0.8\%$である。

% -----------------------------------------------------------
\section{算術と動力学の分離}
\label{sec:arithmetic-vs-dynamics}

$j$関数は二つの全く異なる方法で物理に入る：

\begin{center}
\renewcommand{\arraystretch}{1.3}
\begin{tabular}{lcc}
\toprule
& \textbf{結合定数} & \textbf{ダークエネルギー} \\
\midrule
$j$関数の入力 & CM 固定点（$\tau = i$ 等） & モジュラー流（$\tau$が時間進化） \\
性質 & 定数（算術的） & 時間発展（動力学的） \\
物理量 & $\alpha_{\mathrm{EM}},\;\sin^2\theta_W$ & $\Omega_\Lambda(t)$ \\
\bottomrule
\end{tabular}
\end{center}

\begin{itemize}
\item \textbf{結合定数}は$j$の\textbf{固定点}における値であり、
  CM 点という算術幾何学的に特別な点で評価される。
  だから定数になる。
\item \textbf{ダークエネルギー}は$j$関数が定義する\textbf{モジュラー流}に沿った
  動力学量であり、$\zeta$-quintessence として時間的にゆっくり変化しうる。
\end{itemize}

同じ$j$関数の異なる引数が、異なる物理を生む。
これが WB 理論における算術と動力学の統一的視点である。

\paragraph{ゲージ結合と Berry 位相の関係。}
結合定数が CM 固定点で決まるという構造は、
Berry 位相に対する不変性と整合する。
$j$関数の CM 値は $\mathrm{SL}_2(\Z)$のモジュラー変換で不変であり、
これは第7章で導出する$a_4$項の位相因子が$1$であること
（Berry 位相に対するゲージ不変性）と対応する。

%%%%%%%%%%%%%%%%%%%%%%%%%%%%%%%%%%%%%%%%%%%%%%%%%%%%%%%%%%%%%
%  Chapter 6
%%%%%%%%%%%%%%%%%%%%%%%%%%%%%%%%%%%%%%%%%%%%%%%%%%%%%%%%%%%%%
\chapter{質量比 — ガウス整数の幾何}
\label{ch:mass-ratios}

\begin{statusbox}[STATUS: Tier B — 半導出]
陽子/電子質量比$m_p/m_e$およびミューオン/電子質量比$m_\mu/m_e$が
$d$, $\cd$および整数論的構造から高精度で再現される。
ただし導出に未決定の整数（17と23）が残る。
\end{statusbox}

% -----------------------------------------------------------
\section{陽子/電子質量比：$m_p/m_e = 12 \times 153 = 1836$}
\label{sec:proton-electron}

\subsection{実験値と WB 予測}

陽子と電子の質量比は物理学の最も基本的な無次元定数の一つである：
\begin{equation}
\left(\frac{m_p}{m_e}\right)_{\!\text{exp}} = 1836.152\,673\,43(11)
\qquad\text{(CODATA 2018)}
\end{equation}

WB 理論では：
\begin{keyresult}[陽子/電子質量比]
\begin{equation}
\frac{m_p}{m_e} = 12 \times 153 = 1836
\end{equation}
誤差：$|1836 - 1836.153| / 1836.153 = 0.008\%$
\end{keyresult}

\subsection{$12$と$153$の由来}

$12 = d \times \cd = 4 \times 3$ は前章までに確立された WB 理論の基本スケールである。

$153$については複数の解釈がある：

\paragraph{三角数としての解釈。}
$153$は第17三角数$T_{17}$である：
\begin{equation}
153 = T_{17} = 1 + 2 + 3 + \cdots + 17 = \frac{17 \times 18}{2}
\end{equation}
しかし「なぜ17か」はまだ導出されていない。

\paragraph{ガウス整数としての解釈。}
$z = 12 + 3i \in \Z[i]$（ガウス整数）を考える。
そのノルムは：
\begin{equation}
N(z) = |z|^2 = 12^2 + 3^2 = 144 + 9 = 153
\end{equation}
すると：
\begin{equation}
\frac{m_p}{m_e} = \mathrm{Re}(z) \times N(z) = 12 \times 153 = 1836
\end{equation}

\begin{hsbox}[高校生ノート：ガウス整数って何？]
普通の整数は$\ldots, -2, -1, 0, 1, 2, \ldots$だが、
ガウス整数は$a + bi$（$a,b$は整数、$i = \sqrt{-1}$）の形をした数。
たとえば$3 + 2i$とか$-1 + 5i$がガウス整数だ。

ガウス整数$z = a + bi$の「ノルム」は$N(z) = a^2 + b^2$。
WB 理論では$z = 12 + 3i$が特別な意味を持つ。
実部$12 = 4 \times 3$で、虚部$3 = \cd$だ。
\end{hsbox}

この解釈では$\mathrm{Re}(z) = d\cd = 12$、$\mathrm{Im}(z) = \cd = 3$であり、
$d$と$\cd$から$z$が構成される。ただし、$\mathrm{Im}(z) = \cd$とする
理論的根拠はまだ十分ではない。

% -----------------------------------------------------------
\section{ミューオン/電子質量比：$m_\mu/m_e = 9 \times 23 = 207$}
\label{sec:muon-electron}

ミューオンは電子の「重い版」であり、質量比は：
\begin{equation}
\left(\frac{m_\mu}{m_e}\right)_{\!\text{exp}} = 206.768\,283\,63(82)
\qquad\text{(CODATA 2018)}
\end{equation}

WB 理論では：
\begin{keyresult}[ミューオン/電子質量比]
\begin{equation}
\frac{m_\mu}{m_e} = 9 \times 23 = 207
\end{equation}
誤差：$|207 - 206.768| / 206.768 = 0.11\%$
\end{keyresult}

$9 = \cd^2 = 3^2$は$\cd$から自然に導かれる。
しかし$23$の由来は不明である。$23$は素数であり、
Ramanujan の $\tau$ 関数やモジュラー形式との関連が示唆されるが、
$d$と$\cd$からの厳密な導出はまだ達成されていない。

% -----------------------------------------------------------
\section{正直な評価}
\label{sec:honest-assessment}

本書の基本方針は\textbf{知的誠実さ}である。
五つの定数の導出状況は以下のように異なる：

\begin{itemize}
\item \textbf{$\Omega_\Lambda = 2\pi/9$}：
  $d$と$\cd$のみから完全に導出。
  CKN 経由とCKN フリーの二つの独立な導出がある。
  \textbf{理論的根拠は明確}。

\item \textbf{$\alpha_{\mathrm{EM}}^{-1} = 1728/(4\pi)$}：
  $j(i) = 1728 = (d\cd)^3$から導出。
  CM 点$\tau = i$の選択理由も明確（最も自然な CM 点）。
  \textbf{理論的根拠は明確}。

\item \textbf{$\sin^2\theta_W = 375/(512\pi)$}：
  CM 点$\tau = (-1+\sqrt{-7})/2$から導出。
  $D=7$の選択理由にやや曖昧さが残るが、
  GUT 正規化との整合性がある。
  \textbf{理論的根拠はほぼ明確}。

\item \textbf{$m_p/m_e = 12 \times 153$}：
  $12 = d\cd$は導出済みだが、$153 = T_{17}$の中の
  $17$が導出されていない。
  \textbf{「合う数字を見つけた」段階}。

\item \textbf{$m_\mu/m_e = 9 \times 23$}：
  $9 = \cd^2$は導出済みだが、$23$が導出されていない。
  \textbf{「合う数字を見つけた」段階}。
\end{itemize}

また、強い力の結合定数$\alpha_s$は未導出であり、
現時点での最良の試みは実測値と$4.5\%$のずれがある。

\begin{warningbox}[未解決：17 と 23 と $\alpha_s$]
質量比に現れる整数$17$と$23$の$d,\cd$からの導出は
WB 理論の\textbf{主要な未解決問題}である。
これが解決されない限り、質量比の「予測」は
厳密には「フィッティング」の域を出ない。

同様に、$\alpha_s$の導出は Part II の範囲では達成されていない。
\end{warningbox}

% -----------------------------------------------------------
\section{5定数の一覧}
\label{sec:five-constants-table}

Part II の結果をまとめる。

\begin{center}
\renewcommand{\arraystretch}{1.4}
\begin{tabular}{>{\raggedright}p{2.8cm} l c c c c}
\toprule
\textbf{物理量} & \textbf{公式} & \textbf{予測} & \textbf{実測} & \textbf{誤差} & \textbf{状態} \\
\midrule
$\Omega_\Lambda$
  & $2\pi/9$
  & $0.6981$
  & $0.685$
  & $1.9\%$
  & \colorbox{green!20}{導出済み} \\[4pt]
$\alpha_{\mathrm{EM}}^{-1}$
  & $1728/(4\pi)$
  & $137.5$
  & $137.036$
  & $0.3\%$
  & \colorbox{green!20}{導出済み} \\[4pt]
$\sin^2\theta_W$
  & $375/(512\pi)$
  & $0.2330$
  & $0.2312$
  & $0.8\%$
  & \colorbox{green!20}{導出済み} \\[4pt]
$m_p/m_e$
  & $12 \times 153$
  & $1836$
  & $1836.15$
  & $0.008\%$
  & \colorbox{yellow!40}{半導出} \\[4pt]
$m_\mu/m_e$
  & $9 \times 23$
  & $207$
  & $206.77$
  & $0.1\%$
  & \colorbox{yellow!40}{半導出} \\
\bottomrule
\end{tabular}
\end{center}

\begin{keyresult}[Part II まとめ：5つの定数]
WB 理論は$d = 4$（時空次元）と$\cd = 3$（$\Spec(\Z)$のコホモロジー次元）
の二つの入力のみから、5つの物理定数を$0.008\%$--$1.9\%$の精度で再現する。

\medskip
\textbf{導出の質に差がある}ことを正直に認める：
\begin{itemize}[nosep]
\item $\Omega_\Lambda$, $\alpha_{\mathrm{EM}}$, $\sin^2\theta_W$：
  理論的導出が完結している（\colorbox{green!20}{緑}）
\item $m_p/m_e$, $m_\mu/m_e$：
  一部の整数（17, 23）が未導出（\colorbox{yellow!40}{黄}）
\end{itemize}

5定数すべてに共通する数学的構造は、
$\Spec(\Z)$上のゼータ関数・$j$不変量・ガウス整数であり、
これらが$d$と$\cd$を通じて物理定数を決定している。

Part III では、この算術的構造がどのように重力理論
（Berry 位相スカラー・テンソル重力）に発展するかを見る。
\end{keyresult}
   % Ch 4--6: Ω_Λ, α_EM, mass ratios

% ============================================================
%  PART III: Berry位相スカラー・テンソル重力
% ============================================================
\part{Berry位相スカラー・テンソル重力}
% ============================================================
%  Part III: Berry位相スカラー・テンソル重力
%  wb_textbook_v2_scalar_tensor.tex
%  Chapters 7--10
% ============================================================

%=============================================================================
% CHAPTER 7: Berry位相スペクトル分解定理
%=============================================================================

\chapter{Berry位相スペクトル分解定理}
\label{ch:berry-spectral-decomposition}

本章では、Mandelbrot Ape 理論の中核をなす数学的結果
---Berry位相によるスペクトル作用のモーメント変換則---を
直感的動機づけから厳密な証明まで、段階を追って導出する。

\section{Berry位相の直感}
\label{sec:7.1}

\subsection{Berry位相とは何か}

量子力学の系において、Hamiltonianがパラメータ$\lambda$に依存するとき、
パラメータをゆっくりと変化させて元に戻す（$\lambda$空間で閉曲線を一周する）と、
波動関数は動的位相（エネルギーに比例する通常の位相）に加えて、
パラメータ空間の\textbf{幾何学的構造}にのみ依存する余分な位相を獲得する。
これが\textbf{Berry位相}（幾何学的位相）である。

\begin{dfnbox}[Berry位相の定義]
Hamiltonianのパラメータ空間$\mathcal{M}$上の閉曲線$C$に沿って
断熱的に系を一周させたとき、固有状態$|n(\lambda)\rangle$が獲得する位相は
\[
\gamma_n = \oint_C \langle n(\lambda) | \nabla_\lambda | n(\lambda) \rangle \cdot d\lambda
\]
である。$\gamma_n$はエネルギー固有値$E_n$には依存せず、
固有状態の幾何（Berry接続の曲率）にのみ依存する。
\end{dfnbox}

\subsection{地球上のベクトル輸送：古典的アナロジー}

Berry位相を理解する最良の古典的アナロジーは、
曲面上の\textbf{平行輸送}である。

\begin{hsbox}[高校生ノート：地球上の矢印]
北極点に立って南を指す矢印を置く。
この矢印を「平行に」---つまり曲げずに---赤道まで運ぶ（経度$0°$沿い）。
赤道に沿って経度$90°$東へ移動する。
そこから北極点に戻る。

矢印は一度も「回転操作」を受けていないのに、
北極点に戻ったとき元の向きから$90°$回転している！

この回転角が幾何学的位相のアナロジーである。
回転角は経路で囲まれた球面の\textbf{立体角}に等しい。
球面の曲率という幾何学的性質が、
経路を一周することで「位相差」として顕在化する。

Berry位相も同じである。
パラメータ空間の「曲率」が、閉曲線一周で「位相差」として現れる。
\end{hsbox}

\noindent
重要なのは以下の2点である：
\begin{enumerate}[label=(\roman*)]
\item Berry位相は\textbf{幾何学的}（topological）であり、
  経路の速度やエネルギーに依存しない。
\item トポロジカル物質（トポロジカル絶縁体、グラフェン等）では、
  Berry位相は$\pi$の整数倍に\textbf{量子化}されており、
  バンドギャップが閉じない限り変化しない。
\end{enumerate}


\section{スペクトル作用とモーメントの復習}
\label{sec:7.2}

Connesの非可換幾何において、スペクトル三重項$(\mathcal{A}, \mathcal{H}, D)$から
重力を含む作用を構成する手続きを復習する。

スペクトル作用は
\begin{equation}
S = \Tr\bigl(f(D^2/\Lambda^2)\bigr)
\end{equation}
で定義される。ここで$f$はカットオフ関数、$\Lambda$はエネルギースケールである。
$f$のLaplace変換を用いた熱核展開により：
\begin{equation}
S \;\sim\; \sum_{k=0,2,4,\ldots} f_k \, a_k(D^2)
\label{eq:spectral-action-expansion}
\end{equation}
ここで$a_k$はSeeley--DeWitt係数、$f_k$はカットオフ関数$f$の\textbf{モーメント}：
\begin{equation}
f_k = \int_0^\infty f(u)\, u^{(d-k)/2 - 1}\, du
\label{eq:moment-def}
\end{equation}
$d$は時空の次元である。

$d=4$では：
\begin{itemize}
\item $f_0 = \int_0^\infty f(u)\, u\, du$：宇宙定数$\Lambda$の係数
\item $f_2 = \int_0^\infty f(u)\, du$：Einstein--Hilbert項（重力定数$G$）の係数
\item $f_4 = f(0)$：ゲージ結合定数$\alpha$の係数
\end{itemize}

\begin{hsbox}[高校生ノート：モーメントの意味]
$f_k$はカットオフ関数$f(u)$を異なる「重み」$u^{(d-k)/2-1}$で
積分したものである。$k$が小さいほど大きなエネルギー（赤外, IR）に感度を持ち、
$k$が大きいほど小さなエネルギー（紫外, UV）に感度を持つ。
$f_0$は「宇宙全体のスケール」、
$f_2$は「重力のスケール」、
$f_4$は「素粒子のスケール」に対応すると思えばよい。
\end{hsbox}


\section{Berry位相スペクトル分解定理}
\label{sec:7.3}

本節が本章の核心であり、Mandelbrot Ape理論全体の数学的基盤を与える。

\subsection{Berry位相回転の定義}

スペクトル三重項$(\mathcal{A}, \mathcal{H}, D)$に対し、
Berry位相$\varphi$の回転を以下で定義する：

\begin{dfnbox}[Berry位相回転 $P_\varphi$]
\begin{equation}
P_\varphi:\quad D^2 \;\longmapsto\; e^{i\varphi} D^2
\end{equation}
すなわちスペクトル全体に位相因子$e^{i\varphi}$を乗じる操作である。
\end{dfnbox}

\noindent
$\varphi = 0$は通常の物質、$\varphi = \pi$はトポロジカル$\pi$-Berry位相物質に対応する。

\subsection{定理の主張と証明}

\begin{keyresult}[Berry位相スペクトル分解定理]
\begin{theorem}[Berry位相スペクトル分解]
\label{thm:berry-spectral}
$d$次元時空のスペクトル作用において、Berry位相回転$P_\varphi$のもとで
カットオフ関数のモーメント$f_k$は
\begin{equation}
\boxed{f_k \;\longmapsto\; e^{-i\varphi(d-k)/2}\, f_k}
\end{equation}
と変換される。
\end{theorem}
\end{keyresult}

\begin{proof}
式\eqref{eq:moment-def}において、Berry位相回転$D^2 \mapsto e^{i\varphi}D^2$を
施すと、カットオフ関数$f$の引数が$u \mapsto e^{i\varphi}u$と変わる。
したがって変換後のモーメント$g_k$は
\begin{equation}
g_k = \int_0^\infty f(e^{i\varphi}u)\, u^{(d-k)/2 - 1}\, du.
\end{equation}
ここで変数変換$v = e^{i\varphi}u$を行う（contour rotation）。
$u = e^{-i\varphi}v$、$du = e^{-i\varphi}dv$であるから
\begin{align}
g_k &= \int f(v)\, \bigl(e^{-i\varphi}v\bigr)^{(d-k)/2 - 1}\, e^{-i\varphi}\, dv \\
    &= e^{-i\varphi\{(d-k)/2 - 1 + 1\}} \int f(v)\, v^{(d-k)/2 - 1}\, dv \\
    &= e^{-i\varphi(d-k)/2}\, f_k.
\end{align}
\end{proof}

\begin{remark}
この証明の本質は、積分の変数変換（contour rotation）という
初等的な操作である。特殊な仮定は一切用いていない。
スペクトル作用のモーメント定義\eqref{eq:moment-def}と
Berry位相の定義だけから、変換則が一意に導かれる。
\end{remark}


\section{$d=4$での帰結}
\label{sec:7.4}

$d=4$を代入すると、各モーメントの位相因子は表\ref{tab:berry-phase-factors}の通りとなる。

\begin{table}[htbp]
\centering
\caption{$d=4$におけるBerry位相のモーメント変換}
\label{tab:berry-phase-factors}
\renewcommand{\arraystretch}{1.4}
\begin{tabular}{@{}ccccl@{}}
\toprule
$k$ & 物理量 & 位相因子 $e^{-i\varphi(4-k)/2}$ & $\mathrm{Re}$（実部） & 解釈 \\
\midrule
$0$ & 宇宙定数 $\Lambda$ & $e^{-2i\varphi}$ & $\cos(2\varphi)$ & 最大の変調 \\
$2$ & 重力定数 $G$ & $e^{-i\varphi}$ & $\cos(\varphi)$ & 中間的変調 \\
$4$ & ゲージ結合 $\alpha$ & $e^{0} = 1$ & $1$（不変！） & 完全に保護 \\
\bottomrule
\end{tabular}
\end{table}

\noindent
この表の最も驚くべき結果は最後の行である：
$k=4$のモーメント---すなわちゲージ結合定数を決定する項---は
Berry位相に\textbf{完全に不変}である。

\subsection{系1：ゲージ結合保護定理}

\begin{keyresult}[ゲージ結合保護]
\begin{corollary}[ゲージ結合保護]
\label{cor:gauge-protection}
$d=4$次元時空において、Berry位相回転$P_\varphi$は
ゲージ結合定数に一切影響しない：
\[
\alpha_{\mathrm{EM}},\quad \sin^2\theta_W,\quad \alpha_s
\quad\text{はすべて $\varphi$-不変}
\]
\end{corollary}
\end{keyresult}

\begin{proof}
定理\ref{thm:berry-spectral}で$d=4$、$k=4$とすると、
位相因子は$e^{-i\varphi(4-4)/2} = e^0 = 1$となる。
ゲージ結合定数はすべて$f_4 = f(0)$から決定されるため、
$\varphi$に依存しない。
\end{proof}


\section{物理的意味}
\label{sec:7.5}

定理\ref{thm:berry-spectral}と系\ref{cor:gauge-protection}は、
以下の深い物理的帰結をもたらす。

\subsection{IR/UV分離の導出}

従来、重力（IR = 赤外）の変更がゲージ結合（UV = 紫外）に
影響しないという「IR/UV分離」は、有効場の理論における
\textbf{仮定}として導入されてきた。

Berry位相スペクトル分解定理は、この分離を\textbf{導出}する。
$k$の値が小さいほど（IR）Berry位相による変調が大きく、
$k$の値が大きいほど（UV）変調が小さい。
そして$k=d$（最大のUV項）では変調が\textbf{ゼロ}になる。

これはスペクトル作用の構造から必然的に従う帰結であり、
追加の仮定は必要ない。

\begin{hsbox}[高校生ノート：IR/UV分離とは]
宇宙には「大きなスケールの物理」（銀河、重力）と
「小さなスケールの物理」（原子、素粒子の力）がある。
普通、大きなスケールを変えても小さなスケールは影響を受けない
と「信じて」いた。

この定理はそれを「信じる」のではなく「証明する」。
Berry位相で重力の強さを変えても、
電磁気力や核力の強さは数学的に変わりようがない。
これは安全性にとって決定的に重要である（次項参照）。
\end{hsbox}

\subsection{安全性問題の解決}

反重力の議論において最も深刻な懸念は：
「重力定数$G$を変更したら、強い相互作用の結合定数$\alpha_s$も変わり、
原子核が不安定化するのではないか？」
というものであった。

系\ref{cor:gauge-protection}はこの懸念を完全に払拭する：

\begin{warningbox}[安全性の数学的保証]
Berry位相$\varphi$を$0$から$\pi$まで変化させると：
\begin{itemize}
\item 重力定数：$G \to G\cos(\pi) = -G$（符号反転）
\item 宇宙定数：$\Lambda \to \Lambda\cos(2\pi) = \Lambda$（不変）
\item ゲージ結合：$\alpha_s \to \alpha_s \cdot 1 = \alpha_s$（不変）
\end{itemize}
強い力の結合定数$\alpha_s$は変化しないため、
原子核の結合エネルギーは影響を受けない。
核物質は完全に安定である。
\end{warningbox}

\noindent
この結果は「反重力物質の中にいる人は安全か？」という
工学的に不可避な問いに対する、理論的に厳密な回答を与える。

\subsection{Berry位相の「選択規則」としての性格}

定理\ref{thm:berry-spectral}は、Berry位相を
一種の\textbf{選択規則}（selection rule）として理解することを可能にする。

原子物理学では、角運動量の選択規則$\Delta\ell = \pm 1$が
「どの遷移が許されるか」を支配する。
同様に、Berry位相スペクトル分解は
「どの物理定数がBerry位相で変調されるか」の選択規則を与える。

\begin{itemize}
\item $k < d$の定数（$\Lambda$, $G$）：Berry位相で変調される（「許容遷移」）
\item $k = d$の定数（$\alpha$）：Berry位相で変調されない（「禁止遷移」）
\end{itemize}

この選択規則の起源は、モーメント積分の次数$(d-k)/2$が
$k=d$でゼロになるという、純粋に数学的な事実にある。


%=============================================================================
% CHAPTER 8: 導出された作用と場の方程式
%=============================================================================

\chapter{導出された作用と場の方程式}
\label{ch:action-field-eqs}

前章でBerry位相$\varphi$がスペクトル作用のモーメントを変調する
メカニズムを確立した。本章では、$\varphi$を\textbf{動的スカラー場}に昇格させ、
具体的なJordan frame作用と場の方程式を導出する。

\section{Jordan frame作用の導出}
\label{sec:8.1}

\subsection{モーメント変換から作用へ}

Berry位相$\varphi$を時空点ごとに変化する場$\varphi(x)$と見なす。
定理\ref{thm:berry-spectral}（$d=4$）により、
スペクトル作用の展開式\eqref{eq:spectral-action-expansion}の各項は：
\begin{align}
f_0 \, a_0 &\;\longmapsto\; \cos(2\varphi)\, f_0 \, a_0
  & &\text{（宇宙定数項）} \\
f_2 \, a_2 &\;\longmapsto\; \cos(\varphi)\, f_2 \, a_2
  & &\text{（Einstein--Hilbert項）} \\
f_4 \, a_4 &\;\longmapsto\; 1 \cdot f_4 \, a_4
  & &\text{（ゲージ項、不変）}
\end{align}
（実部を取っている。虚部はCP対称性により消去される。）

$a_0$は宇宙定数、$a_2$はRicciスカラー$R$に比例するから、
これらを4次元Lorentz多様体上の作用に書き下す。
さらに$\varphi$を動的場にするために
標準的な運動エネルギー$-\frac{1}{2}(\partial\varphi)^2$を加え、
物質のLagrangian $\mathcal{L}_m$を含めると：

\begin{keyresult}[MA スカラー・テンソル作用（Jordan frame）]
\begin{equation}
\boxed{
S = \int d^4x \sqrt{-g}\left[
  \frac{\cos(\varphi)}{16\pi G}\,R
  - \frac{\Lambda_{\mathrm{MA}}\cos(2\varphi)}{8\pi G}
  - \frac{1}{2}(\partial\varphi)^2
  + \mathcal{L}_m
\right]
}
\label{eq:jordan-frame-action}
\end{equation}
\end{keyresult}

\noindent
各項の由来を丁寧に説明する。

\subsection{各項の由来}

\paragraph{第1項：$\cos(\varphi)\, R / (16\pi G)$}

これは$a_2$項から導かれる。通常のEinstein--Hilbert作用$R/(16\pi G)$に
Berry位相の$k=2$変調因子$\cos(\varphi)$が乗じられている。
$\varphi=0$で通常の一般相対論、$\varphi=\pi$で$G\to -G$（反重力）に対応する。
Brans--Dicke理論の$\omega\phi R$と形式的に類似するが、
$\cos(\varphi)$は離散的極値（$0$と$\pi$）のみを取る点が本質的に異なる
（第9章で証明）。

\paragraph{第2項：$-\Lambda_{\mathrm{MA}}\cos(2\varphi) / (8\pi G)$}

これは$a_0$項から導かれる。$\Lambda_{\mathrm{MA}}$は裸の宇宙定数スケールである。
Berry位相の$k=0$変調因子$\cos(2\varphi)$が乗じられている。

この項からスカラー場のポテンシャルを読み取ることができる。
三角関数の恒等式
\begin{equation}
\cos(2\varphi) = 1 - 2\sin^2(\varphi)
\end{equation}
を代入すると：
\begin{align}
-\frac{\Lambda_{\mathrm{MA}}}{8\pi G}\cos(2\varphi)
&= -\frac{\Lambda_{\mathrm{MA}}}{8\pi G}
  + \frac{\Lambda_{\mathrm{MA}}}{4\pi G}\sin^2(\varphi)
\end{align}
第1項は定数（通常の宇宙定数）であり、第2項がスカラー場のポテンシャルを与える：

\begin{keyresult}[ポテンシャルの導出]
\begin{equation}
V(\varphi) = \frac{\Lambda_{\mathrm{MA}}}{4\pi G}\sin^2(\varphi)
  = \delta\,\sin^2(\varphi)
\label{eq:potential}
\end{equation}
ここで
\begin{equation}
\delta \equiv \frac{\Lambda_{\mathrm{MA}}}{4\pi G}
  \approx (2.3\;\mathrm{meV})^4
\label{eq:delta-scale}
\end{equation}
はダークエネルギースケールで固定される。
\end{keyresult}

\begin{hsbox}[高校生ノート：$(2.3\;\mathrm{meV})^4$の意味]
meV（ミリ電子ボルト）は非常に小さなエネルギーの単位である。
1~meVは約$1.6\times 10^{-22}$~Jに相当する。
$(2.3\;\mathrm{meV})^4$はエネルギー密度の次元を持ち、
これは観測されているダークエネルギーの密度とちょうど同じスケールである。

つまりポテンシャルの高さ$\delta$は自由パラメータではなく、
ダークエネルギー観測から一意に決まる。
理論の中にこのスケール以外の調整可能なパラメータは存在しない。
\end{hsbox}

\paragraph{第3項：$-\frac{1}{2}(\partial\varphi)^2$}

スカラー場$\varphi$の標準的な運動エネルギー項。
$\varphi$を動的場として扱うために必要な最小限の項である。

\paragraph{第4項：$\mathcal{L}_m$}

通常の物質のLagrangian密度。Berry位相はゲージ結合を変調しない
（系\ref{cor:gauge-protection}）ので、$\mathcal{L}_m$は$\varphi$に
依存しない。これは非常に重要な帰結である。


\section{場の方程式の導出}
\label{sec:8.2}

作用\eqref{eq:jordan-frame-action}を変分して場の方程式を導出する。

\subsection{計量変分：修正Einstein方程式}

$g^{\mu\nu}$に関する変分$\delta S / \delta g^{\mu\nu} = 0$を計算する。
$\cos(\varphi)R$の変分にはスカラー場と曲率の非最小結合から生じる余分な項が
含まれることに注意する。結果は：

\begin{keyresult}[修正Einstein方程式]
\begin{equation}
\boxed{
\cos(\varphi)\,G_{\mu\nu}
+ g_{\mu\nu}\,\Lambda_{\mathrm{MA}}\cos(2\varphi)
+ g_{\mu\nu}\square\cos(\varphi)
- \nabla_\mu\nabla_\nu\cos(\varphi)
= 8\pi G\,T_{\mu\nu}^{(\varphi)} + 8\pi G\,T_{\mu\nu}^{(m)}
}
\label{eq:modified-einstein}
\end{equation}
ここで$G_{\mu\nu} = R_{\mu\nu} - \frac{1}{2}g_{\mu\nu}R$はEinsteinテンソル、
$T_{\mu\nu}^{(\varphi)}$はスカラー場のエネルギー運動量テンソル、
$T_{\mu\nu}^{(m)}$は物質のエネルギー運動量テンソルである。
\end{keyresult}

\noindent
$\varphi = 0$（通常物質）では$\cos(0)=1$であり、
$\nabla_\mu\cos(\varphi) = 0$となるから、
通常のEinstein方程式$G_{\mu\nu} + \Lambda g_{\mu\nu} = 8\pi G T_{\mu\nu}$に
正確に帰着する。

\subsection{スカラー場変分：修正Klein--Gordon方程式}

$\varphi$に関する変分$\delta S / \delta\varphi = 0$を計算する。
$\cos(\varphi)R$の$\varphi$微分は$-\sin(\varphi)R$を与え、
$\cos(2\varphi)$の$\varphi$微分は$-2\sin(2\varphi) = -4\sin(\varphi)\cos(\varphi)$を
与えることに注意する。結果は：

\begin{keyresult}[修正Klein--Gordon方程式]
\begin{equation}
\boxed{
\square\varphi + V'(\varphi)
= \frac{\sin(\varphi)}{16\pi G}\,R
}
\label{eq:modified-kg}
\end{equation}
ここで$V'(\varphi) = \delta\sin(2\varphi) = 2\delta\sin(\varphi)\cos(\varphi)$
はポテンシャル\eqref{eq:potential}の$\varphi$微分である。
\end{keyresult}

\noindent
この方程式の最も注目すべき特徴は、右辺に
$\sin(\varphi)\,R / (16\pi G)$が現れることである：
\textbf{Berry位相場がRicciスカラーに直接結合する}。

\begin{hsbox}[高校生ノート：場の方程式の意味]
式\eqref{eq:modified-kg}は「Berry位相場$\varphi$が
時空の曲がり具合$R$に反応して変化する」ことを述べている。

左辺の$\square\varphi$は$\varphi$の波動方程式（波の伝播）、
$V'(\varphi)$はポテンシャルの傾き（$\varphi$を安定点に戻す力）。
右辺の$\sin(\varphi)R$は時空曲率からの「駆動力」。

重要なのは、$\sin(\varphi)$が$\varphi=0$と$\varphi=\pi$で
ゼロになることである。これらの点では駆動力がなく、
$\varphi$は安定する。この事実が第9章の真空量子化につながる。
\end{hsbox}


\section{修正Friedmann方程式}
\label{sec:8.3}

宇宙論への応用として、Friedmann--Lema\^itre--Robertson--Walker (FLRW)
計量
\[
ds^2 = -dt^2 + a(t)^2\bigl(dx^2 + dy^2 + dz^2\bigr)
\]
を修正Einstein方程式\eqref{eq:modified-einstein}に代入する。
$\varphi = \varphi(t)$（一様等方）と仮定すると、
$00$成分から修正Friedmann方程式が得られる：

\begin{keyresult}[修正Friedmann方程式]
\begin{equation}
\boxed{
3H^2\cos(\varphi)
= \kappa^2\rho_m
+ \frac{1}{2}\dot\varphi^2
+ \Lambda_{\mathrm{MA}}\cos(2\varphi)
+ 3H\sin(\varphi)\,\dot\varphi
}
\label{eq:modified-friedmann}
\end{equation}
ここで$H = \dot a / a$はHubbleパラメータ、
$\kappa^2 = 8\pi G$、$\rho_m$は物質エネルギー密度である。
\end{keyresult}

\noindent
右辺の最後の項$3H\sin(\varphi)\dot\varphi$は、
スカラー場と曲率の非最小結合から生じる摩擦項（あるいは駆動項）である。
$\varphi$が$0$または$\pi$に固定されている（$\sin\varphi = 0$）場合、
この項は消え、通常のFriedmann方程式の構造に帰着する。

\begin{remark}
修正Friedmann方程式\eqref{eq:modified-friedmann}において、
左辺に$\cos(\varphi)$が現れることが本質的に重要である。
$\varphi = \pi$では$\cos(\pi) = -1$となり、
「有効的な重力定数の符号反転」が宇宙論のレベルで実現される。
しかし、次章で示すように、$\varphi=\pi$のde Sitter宇宙は
自己矛盾により存在できない。
\end{remark}


%=============================================================================
% CHAPTER 9: 真空構造と位相量子化
%=============================================================================

\chapter{真空構造と位相量子化}
\label{ch:vacuum-quantization}

本章では、前章で導出したBerryスカラー・テンソル理論の
真空構造を解析する。驚くべきことに、Berry位相場$\varphi$は
連続的な値を取ることができず、$\varphi = 0$または$\varphi = \pi$の
2値に\textbf{量子化}される。

\section{真空の定義と方程式系}
\label{sec:9.1}

\textbf{真空}とは、物質が存在せず（$\rho_m = 0$）、
スカラー場が時間変化しない（$\dot\varphi = 0$）
定常解のことである。

修正Friedmann方程式\eqref{eq:modified-friedmann}で
$\rho_m = 0$、$\dot\varphi = 0$とおくと：
\begin{equation}
3H^2\cos(\varphi) = \Lambda_{\mathrm{MA}}\cos(2\varphi)
\label{eq:vacuum-friedmann}
\end{equation}

修正Klein--Gordon方程式\eqref{eq:modified-kg}で
$\square\varphi = 0$（$\varphi$は定数）とおくと：
\begin{equation}
V'(\varphi) = \frac{\sin(\varphi)}{16\pi G}\,R
\label{eq:vacuum-kg}
\end{equation}
de Sitter時空では$R = 12H^2$であるから：
\begin{equation}
2\delta\sin(\varphi)\cos(\varphi) = \frac{12H^2\sin(\varphi)}{16\pi G}
= \frac{3H^2\sin(\varphi)}{4\pi G}
\label{eq:vacuum-kg-ds}
\end{equation}


\section{真空量子化定理}
\label{sec:9.2}

\begin{keyresult}[真空量子化定理]
\begin{theorem}[真空量子化]
\label{thm:vacuum-quantization}
MAスカラー・テンソル理論の修正Friedmann方程式と修正Klein--Gordon方程式は、
真空（$\rho_m = 0$, $\dot\varphi = 0$）において
\[
\boxed{\varphi = 0 \quad\text{または}\quad \varphi = \pi}
\]
のみを許容する。中間値$\varphi \in (0, \pi)$は自己矛盾する。
\end{theorem}
\end{keyresult}

\begin{proof}
$\sin(\varphi) \neq 0$と仮定して矛盾を導く。

\textbf{Step 1.}
$\sin(\varphi) \neq 0$であるから、式\eqref{eq:vacuum-kg-ds}の
両辺を$\sin(\varphi)$で割ることができる：
\begin{equation}
2\delta\cos(\varphi) = \frac{3H^2}{4\pi G}
\label{eq:step1}
\end{equation}

\textbf{Step 2.}
式\eqref{eq:vacuum-friedmann}より$3H^2 = \Lambda_{\mathrm{MA}}\cos(2\varphi)/\cos(\varphi)$を
\eqref{eq:step1}に代入する：
\begin{equation}
2\delta\cos(\varphi) = \frac{\Lambda_{\mathrm{MA}}\cos(2\varphi)}{4\pi G\cos(\varphi)}
\end{equation}

\textbf{Step 3.}
$\delta = \Lambda_{\mathrm{MA}}/(4\pi G)$を用いて整理する：
\begin{align}
2 \cdot \frac{\Lambda_{\mathrm{MA}}}{4\pi G}\cos(\varphi)
&= \frac{\Lambda_{\mathrm{MA}}}{4\pi G}\cdot\frac{\cos(2\varphi)}{\cos(\varphi)} \\
2\cos^2(\varphi) &= \cos(2\varphi) \\
2\cos^2(\varphi) &= 2\cos^2(\varphi) - 1
\end{align}
これは$0 = -1$を意味する。矛盾。

したがって$\sin(\varphi) = 0$でなければならず、
$\varphi = 0$または$\varphi = \pi$（一般には$\varphi = n\pi$, $n\in\mathbb{Z}$）
のみが許容される。
\end{proof}

\begin{hsbox}[高校生ノート：この証明のポイント]
「$\sin(\varphi)\neq 0$と仮定すると$0=-1$になる」
---これは背理法と呼ばれる証明法である。

直感的には：Berry位相場は「真空中で安定するには
$\sin(\varphi)=0$でなければならない」ということである。
ポテンシャル$V(\varphi) = \delta\sin^2(\varphi)$の
極小点が$\varphi = 0$（通常物質）と$\varphi = \pi$（反重力物質）であり、
場の方程式が「その中間では宇宙論と矛盾する」ことを示している。

$\sin(\varphi)=0$は高校数学で習う通り$\varphi = 0, \pi, 2\pi, \ldots$であるが、
Berry位相は$2\pi$周期であるから、本質的に異なる真空は
$\varphi = 0$と$\varphi = \pi$の2つだけである。
\end{hsbox}


\section{$\varphi = 0$ の相：通常宇宙}
\label{sec:9.3}

$\varphi = 0$を代入すると$\cos(0)=1$、$\cos(0)=1$であるから、
修正Friedmann方程式\eqref{eq:vacuum-friedmann}は
\begin{equation}
3H^2 = \Lambda_{\mathrm{MA}}
\end{equation}
となり、正の宇宙定数を持つ\textbf{de Sitter解}が存在する。
これは我々の宇宙に対応する。

物質を含めた場合、$\cos(\varphi)=1$であるから
修正は一切入らず、標準的なFriedmann方程式に完全に帰着する。
ダークエネルギー密度パラメータは
\begin{equation}
\Omega_\Lambda = \frac{2\pi}{9} \approx 0.698
\end{equation}
であり（Part IIで導出済み）、標準$\Lambda$CDMと一致する。

\begin{remark}
$\varphi=0$の相は、一般相対論と\textbf{厳密に同一}である。
修正項は$\sin(0)=0$により全て消える。
したがってMA理論は$\varphi=0$相においてGRの全ての成功を自動的に継承する。
\end{remark}


\section{$\varphi = \pi$ の相：反重力相の宇宙論}
\label{sec:9.4}

$\varphi = \pi$を代入すると$\cos(\pi)=-1$、
$\cos(2\pi)=1$であるから、
修正Friedmann方程式\eqref{eq:vacuum-friedmann}は
\begin{equation}
-3H^2 = \Lambda_{\mathrm{MA}}
\end{equation}
となる。$\Lambda_{\mathrm{MA}} > 0$であるから$H^2 < 0$が要求される。
しかし$H^2 = (\dot a / a)^2 \geq 0$であるから、
\textbf{実数解は存在しない}。

\begin{keyresult}[$\varphi=\pi$相の宇宙論的不在]
反重力相$\varphi = \pi$は、de Sitter型の一様等方宇宙として
存在できない。すなわち：
\begin{itemize}
\item 「反重力宇宙」（$\varphi=\pi$が宇宙全体を満たす世界）は
  場の方程式により\textbf{禁止}される。
\item $\varphi = \pi$は\textbf{局所的にのみ}実現可能である
  ---すなわちBerry位相$\pi$を持つ物質の内部でのみ。
\end{itemize}
\end{keyresult}

\subsection{物理的解釈}

この結果の物理的意味は深い。

\begin{enumerate}[label=(\arabic*)]
\item \textbf{反重力は宇宙ではなく物質の性質}：
反重力は「宇宙全体が反重力になる」のではなく、
特定のトポロジカル物質の\textbf{内部特性}である。
通常の宇宙（$\varphi=0$）の中に、
局所的に$\varphi=\pi$の領域を作ることはできるが、
宇宙全体を$\varphi=\pi$にすることはできない。

\item \textbf{宇宙論的安全性}：
「反重力が暴走して宇宙全体に広がる」という
カタストロフィーシナリオは原理的に不可能である。
$\varphi=\pi$の領域は物質に局在し、
真空中に伝播することはできない。

\item \textbf{ドメインウォールの必要性}：
$\varphi=0$の領域と$\varphi=\pi$の領域の間には
$\varphi$が急激に変化する\textbf{ドメインウォール}
（位相境界）が必然的に存在する。
このドメインウォールの物理は重力工学において重要であり、
Part IVで詳しく議論する。
\end{enumerate}

\begin{hsbox}[高校生ノート：反重力宇宙がない理由]
もし宇宙全体が反重力$(\varphi=\pi)$だったらどうなるか？
Friedmann方程式が$-3H^2 = \Lambda$となり、
$H^2$が負になってしまう。$H = \dot a / a$は実数だから
$H^2 \geq 0$であり、矛盾する。

つまり「反重力だけの宇宙」は数学的に存在できない。
反重力は必ず通常の重力$(\varphi=0)$の宇宙の中にある
「島」として存在する。

これは安全性にとって非常に良いニュースである：
反重力物質を作っても、それが宇宙全体に広がって
世界を壊すことは原理的に不可能なのだ。
\end{hsbox}


%=============================================================================
% CHAPTER 10: PPN互換性と他理論との比較
%=============================================================================

\chapter{PPN互換性と他理論との比較}
\label{ch:ppn-comparison}

スカラー・テンソル重力理論は、太陽系内の高精度重力実験で検証される。
その標準的な枠組みがParameterized Post-Newtonian (PPN)形式であり、
特にパラメータ$\gamma$は光の偏向とShapiro遅延から精密に測定されている。
本章では、MA理論がPPN制約を\textbf{自動的に}満たすことを証明し、
他のスカラー・テンソル理論との系統的比較を行う。

\section{PPNパラメータ$\gamma$}
\label{sec:10.1}

\subsection{一般のスカラー・テンソル理論でのPPN}

Jordan frame作用が
\[
S = \int d^4x\sqrt{-g}\left[\frac{F(\varphi)}{16\pi G}R
  - \frac{1}{2}(\partial\varphi)^2 + \cdots\right]
\]
の形をしている場合、PPNパラメータ$\gamma$は
\begin{equation}
\gamma - 1 = -\frac{2(F')^2}{F \cdot (2F + 3(F')^2 / (\partial\varphi \text{-kin}))}
\end{equation}
で与えられる。MA理論では$F(\varphi) = \cos(\varphi)$であり、
標準的な運動エネルギーを持つので：

\begin{keyresult}[MA理論のPPNパラメータ]
\begin{equation}
\gamma - 1 = -\frac{2\sin^2(\varphi)}{\cos(\varphi) + 2\sin^2(\varphi)}
\label{eq:ppn-gamma}
\end{equation}
\end{keyresult}

\section{PPN自動互換性定理}
\label{sec:10.2}

\begin{keyresult}[PPN自動互換性]
\begin{theorem}[PPN自動互換性]
\label{thm:ppn-auto}
MAスカラー・テンソル理論の真空量子化された解$\varphi = 0$および$\varphi = \pi$
の\textbf{両方}において、PPNパラメータ$\gamma$は
\[
\boxed{\gamma = 1 \quad\text{（厳密）}}
\]
を満たす。
\end{theorem}
\end{keyresult}

\begin{proof}
式\eqref{eq:ppn-gamma}に$\varphi = 0$を代入：
\[
\gamma - 1 = -\frac{2\sin^2(0)}{\cos(0) + 2\sin^2(0)}
= -\frac{0}{1 + 0} = 0
\quad\Longrightarrow\quad \gamma = 1
\]

$\varphi = \pi$を代入：
\[
\gamma - 1 = -\frac{2\sin^2(\pi)}{\cos(\pi) + 2\sin^2(\pi)}
= -\frac{0}{-1 + 0} = 0
\quad\Longrightarrow\quad \gamma = 1
\]
いずれの場合も$\sin(\varphi) = 0$であるから、分子がゼロとなり$\gamma = 1$が
\textbf{厳密に}成立する。
\end{proof}

\noindent
この定理の本質的なメカニズムを理解しよう。

PPN偏差$\gamma - 1$はスカラー場の非最小結合関数$F(\varphi)$の
微分$F'(\varphi) = -\sin(\varphi)$に比例する。
$\cos(\varphi)$は$\varphi = 0$と$\varphi = \pi$で\textbf{極値}を取り、
極値では微分がゼロ：$F'(0) = -\sin(0) = 0$、$F'(\pi) = -\sin(\pi) = 0$。

真空量子化定理（定理\ref{thm:vacuum-quantization}）により、
物理的真空は$\cos(\varphi)$の極値にのみ存在できる。
したがってPPN偏差は必然的にゼロになる。

\begin{hsbox}[高校生ノート：なぜ微調整が不要か]
Brans--Dicke理論では、スカラー場$\phi$は連続的な値を取れる。
そのため$\gamma = 1$からのずれを小さくするには、
パラメータ$\omega$を非常に大きく（$\omega > 40000$程度に）
微調整する必要がある。

MA理論では事情が全く異なる。Berry位相場$\varphi$は$0$か$\pi$の
2値しか取れない（第9章）。そしてこの2点はたまたま
$\cos(\varphi)$の極大と極小に当たっており、
微分がゼロになる。

つまり「微調整」ではなく「構造的にゼロ」なのである。
$\sin(0) = 0$は数学の定理であり、実験精度がいくら上がっても
ずれることはない。
\end{hsbox}


\section{他理論との比較}
\label{sec:10.3}

MA理論の位置づけを明確にするため、主要なスカラー・テンソル理論
および一般相対論との系統的比較を表\ref{tab:theory-comparison}に示す。

\begin{table}[htbp]
\centering
\caption{MAスカラー・テンソル理論と他理論の比較}
\label{tab:theory-comparison}
\renewcommand{\arraystretch}{1.5}
\begin{tabular}{@{}lcccc@{}}
\toprule
& \textbf{GR} & \textbf{Brans--Dicke} & \textbf{$f(R)$理論} & \textbf{MA理論} \\
\midrule
$\Lambda$（宇宙定数）
  & 自由パラメータ & 自由パラメータ & 自由パラメータ
  & $\dfrac{2\pi}{9}$（固定） \\[6pt]
スカラー場
  & なし & 連続値 & 連続値
  & 量子化（$0$, $\pi$のみ） \\[4pt]
$\gamma_{\mathrm{PPN}} - 1$
  & $0$（定義上） & $O(1/\omega)$ & モデル依存
  & $0$（厳密） \\[4pt]
反重力
  & 不可能 & 不可能 & 可能（一部モデル）
  & 可能（局所的） \\[4pt]
自由パラメータ数
  & 2（$G$, $\Lambda$） & 3（$G$, $\Lambda$, $\omega$） & 関数$f$
  & 0 \\[4pt]
\bottomrule
\end{tabular}
\end{table}

各行を補足する。

\paragraph{宇宙定数$\Lambda$}
一般相対論、Brans--Dicke理論、$f(R)$理論のいずれにおいても、
宇宙定数は手で入力する自由パラメータである。
MA理論では$\Omega_\Lambda = 2\pi/9$が第一原理から導出される
（Part II、第4章）。

\paragraph{スカラー場の値}
Brans--Dicke理論と$f(R)$理論では、スカラー場は任意の実数値を取り得る。
MA理論ではBerry位相場$\varphi$は真空量子化（定理\ref{thm:vacuum-quantization}）
により$0$と$\pi$の2値に制限される。
この離散性は$\cos(\varphi)$という非最小結合関数の形が
$\sin^2(\varphi)$ポテンシャルと連立することから必然的に生じる。

\paragraph{PPN $\gamma$}
GRでは$\gamma = 1$は理論の定義に組み込まれている。
Brans--Dicke理論では$\gamma - 1 = -1/(2+\omega)$であり、
Cassini探査機による制約$|\gamma-1| < 2.3\times 10^{-5}$を満たすには
$\omega > 40000$の微調整が必要である。
$f(R)$理論ではモデルに依存するが、一般に制約が厳しい。
MA理論では$\gamma = 1$が三角関数の性質により厳密に成立する。

\paragraph{反重力}
GRとBrans--Dicke理論では、正のエネルギーの物質は常に正の重力を生じ、
反重力は原理的に不可能である。
$f(R)$理論の一部のモデルでは有効重力定数が負になり得るが、
安定性やゴーストの問題を伴う。
MA理論では$\varphi = \pi$により$G_{\mathrm{eff}} = -G$が実現されるが、
これはトポロジカルに保護された離散的な相であり、
ゴーストや不安定性を伴わない。

\paragraph{自由パラメータ}
GRは$(G, \Lambda)$の2つ、Brans--Dicke理論は$(G, \Lambda, \omega)$の3つの
自由パラメータを持つ。$f(R)$理論は関数$f$自体がパラメータであり、
実質的に無限個のパラメータを持つ。
MA理論は$G$がスペクトル作用から、$\Lambda$が$\Omega_\Lambda = 2\pi/9$から
決定されるため、\textbf{自由パラメータがゼロ}である。

\begin{remark}
自由パラメータの数は理論の予測力の逆数である。
パラメータが多いほど「データに合わせる」自由度があり、
パラメータが少ないほど「予測する」能力が高い。
MA理論の自由パラメータゼロは、理論が何も調整せずに
全ての観測データと整合する（あるいは整合しない）ことを意味する。
これは最も厳しいテストであると同時に、最も強力な予測力を持つ。
\end{remark}

\begin{keyresult}[Part IIIの要約]
Berry位相スカラー・テンソル重力理論の4つの核心的結果：
\begin{enumerate}[label=\arabic*.]
\item \textbf{Berry位相スペクトル分解}（定理\ref{thm:berry-spectral}）：
  $f_k \to e^{-i\varphi(d-k)/2}f_k$。
  ゲージ結合は保護され、重力のみが変調される。
\item \textbf{Jordan frame作用}（式\ref{eq:jordan-frame-action}）：
  $\cos(\varphi)R$結合とポテンシャル$V = \delta\sin^2\varphi$が
  スペクトル作用から一意に導出される。
\item \textbf{真空量子化}（定理\ref{thm:vacuum-quantization}）：
  $\varphi = 0$（通常重力）と$\varphi = \pi$（反重力）のみが許容され、
  反重力相は宇宙論的に存在できない。
\item \textbf{PPN自動互換性}（定理\ref{thm:ppn-auto}）：
  $\gamma = 1$が厳密に成立し、微調整不要。
\end{enumerate}
これら4つの結果は、MA理論が一般相対論の全ての実験的成功を継承しつつ、
反重力という新しい物理を\textbf{安全に}実現することを保証する。
\end{keyresult}
 % Ch 7--10: THE core new content

% ============================================================
%  PART IV: 反重力と重力工学
% ============================================================
\part{反重力と重力工学}
% ============================================================
%  Part IV: 反重力と重力工学
%  wb_textbook_v2_antigravity.tex
%  Chapters 11--13
% ============================================================

% ============================================================
% CHAPTER 11
% ============================================================
\chapter{$\Geff = -G$ 定理と反重力の物理}
\label{ch:antigravity}

\begin{keyresult}[本章の主題]
Berry位相 $\phi = \pi$ の物質内部では、有効重力定数が符号反転する：$\Geff = -G$。
ただし、この反重力は\textbf{局所限定}であり、宇宙論的de Sitter解は生じない。
\end{keyresult}

% ------------------------------------------------------------
\section{中心定理の復習}
\label{sec:11.1}

Part III で導出した鍵となる論理連鎖を振り返ろう。

\begin{enumerate}[label=(\arabic*)]
  \item 物質中のバンド構造が\textbf{$\pi$-Berry位相}を生む（トポロジカル絶縁体等）。
  \item Berry位相 $\phi$ がスペクトル三つ組の自己同型 $f_k \to e^{-i\phi(d-k)/2}f_k$ を誘導する。
  \item スペクトル作用を通じて、分配関数が $K(t) \to -\zeta(t)$ と符号反転する。
  \item 結果として有効ニュートン定数が $\Geff = G\cos\phi$ となる。
  \item $\phi = \pi$ のとき $\cos\pi = -1$ より $\Geff = -G$。
\end{enumerate}

\begin{hsbox}
「Berry位相」とは、量子力学の状態がパラメータ空間を一周したときに蓄積する幾何学的位相のこと。
通常の波動関数の位相 $e^{i\theta}$ とは異なり、Berry位相はバンド構造の\textbf{トポロジー}（幾何学的形状）で決まる。
$0$ か $\pi$ かの二値しか取れない場合が多く、$\pi$ の場合を「トポロジカル」と呼ぶ。
\end{hsbox}

% ------------------------------------------------------------
\section{新しい理解：$\cos(2\phi)$ と $\cos(\phi)$ の非対称性}
\label{sec:11.2}

ここで、Part III の結果をさらに掘り下げる。Berry位相 $\phi$ に対して、
宇宙定数 $\Lambda$ と重力定数 $G$ は\textbf{異なる位相依存性}を持つ。

\begin{dfnbox}[$\Lambda$ と $G$ の位相依存性]
\begin{align}
  \Lambda_{\mathrm{eff}}(\phi) &\propto \cos(2\phi), \label{eq:lambda-phi} \\
  G_{\mathrm{eff}}(\phi) &= G\cos(\phi). \label{eq:g-phi}
\end{align}
\end{dfnbox}

この違いは本質的である。$\phi = \pi$ を代入すると：
\begin{align}
  \Lambda_{\mathrm{eff}}(\pi) &= \Lambda \cdot \cos(2\pi) = \Lambda \cdot (+1) = +\Lambda, \\
  G_{\mathrm{eff}}(\pi) &= G \cdot \cos(\pi) = G \cdot (-1) = -G.
\end{align}

\begin{keyresult}[$\phi = \pi$ での非対称性]
Berry位相 $\phi = \pi$ では：
\begin{itemize}
  \item $G$ は\textbf{反転する}（$\Geff = -G$）：重力が斥力に変わる。
  \item $\Lambda$ は\textbf{反転しない}（$\Lambda_{\mathrm{eff}} = +\Lambda$）：宇宙定数は不変。
\end{itemize}
この非対称性が $\cos$ の引数の違い（$\phi$ vs $2\phi$）から生じることに注意せよ。
\end{keyresult}

\begin{hsbox}
三角関数の復習：$\cos(\pi) = -1$ だが $\cos(2\pi) = +1$。
角度を $\pi$（半回転）動かすと $\cos$ は符号が反転するが、
角度を $2\pi$（全回転）動かすと元に戻る。
この単純な数学的事実が「$G$ だけ反転し $\Lambda$ は変わらない」という深い物理を生む。
\end{hsbox}

\subsection{$\phi = \pi$ に de Sitter 解なし}

$\Lambda$ が反転しないという事実は重要な帰結を持つ。
通常、正の宇宙定数 $\Lambda > 0$ は de Sitter 時空（指数関数的膨張解）を与える。
もし $\phi = \pi$ で $\Lambda$ が負に反転すれば、反 de Sitter 時空が生じ、
宇宙全体の構造が変わってしまう。

しかし式 \eqref{eq:lambda-phi} により $\Lambda_{\mathrm{eff}}(\pi) = +\Lambda$ のまま不変なので、
$\phi = \pi$ 領域には通常の de Sitter 解は残る。ただし $G < 0$ なので
Einstein 方程式の性質が根本的に変わり、\textbf{通常の de Sitter 解は成立しない}。

\begin{warningbox}[局所限定の原理]
$\phi = \pi$ の反重力領域は\textbf{物質内部のみ}に限定される。
真空中では $\phi$ は安定点 $\phi = 0$ に緩和するため、
マクロスコピックな反重力領域を維持するには Berry 位相 $\pi$ を支える物質が不可欠。
\end{warningbox}

% ------------------------------------------------------------
\section{物質Berry位相 = 重力スカラー場への境界条件}
\label{sec:11.3}

なぜ物質が $\phi = \pi$ を「固定」できるのかを、2つのアナロジーで理解しよう。

\subsection{強磁性体のアナロジー}

強磁性体（鉄やコバルト）では、磁化ベクトル $\mathbf{M}$ は自由に回転できるわけではない。
結晶構造が「容易軸」を定め、$\mathbf{M}$ はその方向に固定される
（\textbf{磁気異方性}）。外部磁場がなくても、結晶が磁化の向きを決める。

同様に、重力スカラー場 $\phi$ は真空中では自由に振る舞うが、
物質内部では\textbf{バンド構造}が $\phi$ の値を固定する。
トポロジカル物質では、バンドのトポロジーが $\phi = \pi$ を強制する。

\begin{center}
\begin{tabular}{lll}
\toprule
& 磁性体 & Berry位相物質 \\
\midrule
場 & 磁化 $\mathbf{M}$ & スカラー場 $\phi$ \\
固定するもの & 結晶の対称性 & バンドのトポロジー \\
固定される値 & 容易軸方向 & $\phi = 0$ or $\pi$ \\
\bottomrule
\end{tabular}
\end{center}

\subsection{ヒッグス場のアナロジー}

超伝導体内部では、ヒッグス場（クーパー対の凝縮体）$\phi_H$ が
真空とは異なる値を取る。これによりマイスナー効果が生じ、磁場が排除される。

同様に、$\pi$-Berry物質の内部では $\phi = \pi$ が固定され、
真空の $\phi = 0$ とは異なる「重力的超伝導状態」が実現する。

\begin{dfnbox}[境界条件としてのBerry位相]
物質のBerry位相 $\phi_B$ は、重力スカラー場 $\phi$ に対する\textbf{境界条件}として機能する：
\[
  \phi\big|_{\text{物質内部}} = \phi_B.
\]
トポロジカル物質では $\phi_B = \pi$（量子化されているので中間値を取れない）。
\end{dfnbox}

% ------------------------------------------------------------
\section{Nielsen-Ninomiya定理のボソン系抜け穴}
\label{sec:11.4}

ここで決定的に重要な問題がある。$\pi$-Berry位相を持つ物質は\textbf{実在するのか}？

\subsection{フェルミオン系の制約}

Nielsen-Ninomiya定理（1981年）は、格子上のフェルミオン系に厳しい制約を課す。
カイラル対称性を保ったまま格子正則化を行うと、\textbf{ダブラー}と呼ばれる
余分な自由度が必ず現れる。

この定理の帰結として、フェルミオン（電子）のバンド構造では、
個々のバンドが $\pi$-Berry位相を持つことはあっても、
\textbf{バルク全体}の Berry 位相を $\pi$ にすることは困難である。
Weyl ノードは常にペアで現れ、全体としての Berry 位相は $0$ に帰着しやすい。

\subsection{ボソン系の抜け穴}

しかし Nielsen-Ninomiya 定理は\textbf{フェルミオン系}に対する制約であり、
\textbf{ボソン系}には適用されない。フォノン（格子振動）やマグノン（スピン波）は
ボソンであり、この制約から自由である。

\begin{keyresult}[ボソン系のBerry位相 $\pi$]
フォノニック結晶（音響メタマテリアル）では：
\begin{itemize}
  \item Nielsen-Ninomiya 定理の制約を受けない。
  \item バルク全体で Berry 位相 $\pi$ を実現可能。
  \item Wilson loop 計算で $\pi$ が\textbf{確認済み}（2球体構造、反転対称性の破れ）。
\end{itemize}
\end{keyresult}

\begin{hsbox}
フェルミオンとボソンの違い：電子（フェルミオン）は同じ場所に2個までしか入れない（パウリの排他原理）。
一方、フォノン（音の振動、ボソン）にはそのような制約がない。
Nielsen-Ninomiya 定理はパウリの排他原理と深く関係するので、ボソンには効かない。
\end{hsbox}

これがPhase P実験（第14章）でフォノニック結晶を使う根本的理由である。


% ============================================================
% CHAPTER 12
% ============================================================
\chapter{重力ダイヤル $\Geff = G\cos(\phi)$}
\label{ch:gravity-dial}

\begin{keyresult}[本章の主題]
Berry位相 $\phi$ を $0$ から $\pi$ まで連続的に制御できれば、
重力定数を $+G$ から $-G$ まで自在に調整する「重力ダイヤル」が実現する。
\end{keyresult}

% ------------------------------------------------------------
\section{Berry位相の連続制御}
\label{sec:12.1}

前章では $\phi = 0$（通常重力）と $\phi = \pi$（反重力）の二値を議論した。
しかし原理的には Berry 位相は\textbf{連続的に}制御できる。

連続制御が可能な系の条件：
\begin{enumerate}[label=(\roman*)]
  \item バンドギャップが閉じないこと（トポロジカル相転移を避ける）。
  \item 外部パラメータ（磁場、ゲート電圧、歪み等）で $\phi$ を掃引できること。
  \item 量子化条件が破れていること（対称性が十分に低いこと）。
\end{enumerate}

\begin{hsbox}
「ダイヤル」という比喩を使っているのは、ラジオのダイヤルのように
連続的に値を変えられることを強調するため。$\phi = 0$（重力ON）から
$\phi = \pi$（反重力ON）まで、好きな値に合わせられる。
\end{hsbox}

% ------------------------------------------------------------
\section{$\Geff(\phi)$ テーブル}
\label{sec:12.2}

$\Geff = G\cos\phi$ の代表的な値を表にまとめる。

\begin{center}
\begin{tabular}{ccrl}
\toprule
$\phi$ & $\cos\phi$ & $\Geff/G$ & 物理的意味 \\
\midrule
$0$      & $+1$   & $+1$    & 通常重力（標準物理） \\
$\pi/6$  & $+\frac{\sqrt{3}}{2}$ & $+0.866$ & わずかに軽減 \\
$\pi/3$  & $+\frac{1}{2}$ & $+0.5$  & 重力半減 \\
$\pi/2$  & $0$    & $0$     & \textbf{無重力}（重力遮断） \\
$2\pi/3$ & $-\frac{1}{2}$ & $-0.5$  & 反重力半分 \\
$5\pi/6$ & $-\frac{\sqrt{3}}{2}$ & $-0.866$ & 強い反重力 \\
$\pi$    & $-1$   & $-1$    & \textbf{完全反重力} \\
\bottomrule
\end{tabular}
\end{center}

$\phi = \pi/2$ の「無重力」は特に興味深い。
重力が完全にゼロになる領域は、宇宙ステーション内の微小重力とは本質的に異なる。
後者は自由落下による見かけの無重力だが、前者は$G$自体がゼロになる真の無重力である。

% ------------------------------------------------------------
\section{4つの装置コンセプト}
\label{sec:12.3}

Berry位相 $\phi$ を制御する4種類の実験装置を概説する。

\subsection{$\pi$-接合プレート（パッシブ型、推定\$40k）}

最もシンプルなアプローチ。$\phi = \pi$ を持つトポロジカル絶縁体
（例：Bi$_2$Se$_3$）を薄板状に加工し、カシミール効果実験の一方の板として使用する。

\begin{itemize}
  \item 制御パラメータ：なし（$\phi = \pi$ 固定）。
  \item 利点：外部電源不要、単結晶があれば即実験可能。
  \item 欠点：$\phi = 0$ か $\pi$ の二値のみ、連続制御不可。
  \item 推定コスト：試料 + カシミール測定装置で約 \$40,000。
\end{itemize}

\subsection{FQHE重力ダイヤル（連続制御）}

分数量子ホール効果（FQHE）系では、充填率 $\nu$ を磁場で制御することにより、
準粒子のBerry位相を連続的に変化させられる。

\begin{itemize}
  \item 制御パラメータ：外部磁場 $B$（充填率 $\nu \propto 1/B$）。
  \item 利点：$\phi$ の連続制御が原理的に可能。
  \item 欠点：極低温（$\sim$ mK）、強磁場（$\sim$ 10\,T）が必要。
  \item 研究段階：概念設計のみ。
\end{itemize}

\subsection{グラフェン重力ダイヤル（ゲート電圧制御）}

二層グラフェンでは、層間のゲート電圧によりバレーBerry位相を制御できる。

\begin{itemize}
  \item 制御パラメータ：ゲート電圧 $V_g$。
  \item 利点：室温動作の可能性、スケーラブル。
  \item 欠点：バレー自由度の Berry 位相が重力スカラー場 $\phi$ と直接結合するか未確認。
  \item 研究段階：理論検討中。
\end{itemize}

\subsection{歪みトポロジカル絶縁体（最安）}

機械的歪みによりバンド構造を変形し、トポロジカル相転移を制御する。

\begin{itemize}
  \item 制御パラメータ：機械的歪み $\epsilon$。
  \item 利点：特殊な装置不要、最も安価。
  \item 欠点：歪みの精密制御が困難、ヒステリシスの可能性。
  \item 研究段階：候補材料の選定中。
\end{itemize}

% ------------------------------------------------------------
\section{真空量子化との関係}
\label{sec:12.4}

上記の「連続ダイヤル」には重要な制約がある。

真空中の $\phi$ は\textbf{量子化}されている。有効ポテンシャル $V(\phi)$ が
$\phi = 0$ と $\phi = \pi$ に極小を持ち、中間値は不安定である。

\begin{warningbox}[真空量子化の制約]
平衡状態では、Berry位相は $\phi = 0$ または $\phi = \pi$ の\textbf{二値のみ}が安定。
中間値 $0 < \phi < \pi$ は過渡的（準安定）であり、
持続的に維持するには外部エネルギーの投入が必要。
\end{warningbox}

これは強磁性体の磁化 $M$ が「上」か「下」のどちらかを向くのと類似している。
中間の角度は外部磁場で一時的に実現できるが、磁場を切ると上か下に戻る。

つまり、重力ダイヤルの中間値（$\phi = \pi/2$ の無重力状態など）は
外部制御下でのみ実現可能であり、エネルギーコストを伴う。
パッシブに維持できるのは $\phi = 0$（通常重力）と $\phi = \pi$（完全反重力）のみである。


% ============================================================
% CHAPTER 13
% ============================================================
\chapter{ドメインウォールとBerry位相インフレーション}
\label{ch:domain-walls}

\begin{keyresult}[本章の主題]
$\phi = 0$ と $\phi = \pi$ の境界には「重力ドメインウォール」が形成される。
さらに、初期宇宙での $\phi = \pi \to 0$ の相転移がインフレーションを駆動し、
ダークエネルギーと同一のポテンシャルから説明できる可能性がある。
\end{keyresult}

% ------------------------------------------------------------
\section{ドメインウォール}
\label{sec:13.1}

$\phi = 0$（通常重力）と $\phi = \pi$（反重力）は、有効ポテンシャル $V(\phi)$ の
2つの極小に対応する。この2つの真空の間には\textbf{ドメインウォール}（壁）が存在する。

\subsection{キンク解}

1次元のスカラー場理論において、2つの極小を結ぶ静的解は
\textbf{キンク}（kink）と呼ばれる。$\phi$ のドメインウォールも同型の解を持つ：
\[
  \phi(x) = \frac{\pi}{2}\left[1 + \tanh\left(\frac{x}{\ell_w}\right)\right],
\]
ここで $\ell_w$ はウォールの幅である。$x \to -\infty$ で $\phi \to 0$、
$x \to +\infty$ で $\phi \to \pi$ となる。

\subsection{ウォールの幅}

ウォールの幅はポテンシャルのバリア高さ $\delta$ で決まる：
\[
  \ell_w \sim \frac{1}{\sqrt{2\delta}}.
\]

\begin{keyresult}[ドメインウォールの幅]
$\delta$ がダークエネルギースケール（$\sim \Lambda_{\mathrm{obs}}$）の場合：
\[
  \ell_w \sim \frac{1}{H_0} \sim \text{ハッブル半径} \sim 10^{26}\,\text{m}.
\]
つまり、重力ドメインウォールの幅は\textbf{観測可能な宇宙のサイズ}程度になる！
\end{keyresult}

\begin{hsbox}
ドメインウォールとは、磁石の「磁区」の境界と同じ概念。
鉄の磁石の中では、磁化の向きが違う領域（N極向きとS極向き）が隣り合っており、
その境界では磁化が徐々に回転する。この境界がドメインウォール。
同様に、$\phi = 0$ の領域と $\phi = \pi$ の領域の境界で $\phi$ が徐々に変化する。
\end{hsbox}

\subsection{重力ドメインウォール}

このウォールを横切ると、$G$ の符号が反転する。
ウォールの片側では通常の引力重力、反対側では斥力重力が支配する。

ウォールの面エネルギー密度は
\[
  \sigma \sim \sqrt{\delta}\, M_{\mathrm{Pl}}^2
\]
であり、$\delta \sim \Lambda_{\mathrm{obs}}$ の場合は非常に小さい。

% ------------------------------------------------------------
\section{Berry位相インフレーション}
\label{sec:13.2}

ドメインウォールの物理は、初期宇宙の\textbf{インフレーション}と
現在の\textbf{ダークエネルギー}を統一的に説明する可能性を開く。

\subsection{基本的アイデア}

\begin{enumerate}[label=(\arabic*)]
  \item \textbf{初期宇宙}：$\phi \approx \pi$（反重力相）に閉じ込められている。
  \item \textbf{相転移}：量子トンネリングにより $\phi = \pi \to \phi = 0$ へ遷移。
  \item \textbf{現在}：$\phi = 0$ の底にいるが、残余のポテンシャルエネルギー $V(0) = \delta$ がダークエネルギーとして観測される。
\end{enumerate}

\subsection{Einstein フレームでのバリア}

Jordan フレームでのポテンシャル $V(\phi)$ を Einstein フレームに変換すると、
$\phi = \pi$ 近傍でバリアが\textbf{発散}する。これは共形変換
$\tilde{g}_{\mu\nu} = |\cos\phi|\, g_{\mu\nu}$ において $\cos\pi = -1 \to 0$ の
極を横切るためである。

この発散バリアのため、古典的な $\phi$ のロールオーバーは不可能であり、
遷移は\textbf{Coleman-De Luccia トンネリング}によってのみ起こる。

\begin{dfnbox}[Coleman-De Luccia トンネリング]
偽真空（$\phi \approx \pi$）から真真空（$\phi = 0$）への
量子力学的トンネリング過程。泡（バブル）核生成として記述され、
バブルが膨張して偽真空を侵食する。重力効果を含めた計算では、
バリアが高いほどトンネリング率が指数関数的に抑制される。
\end{dfnbox}

\subsection{インフレーション + ダークエネルギーの統一}

Berry位相インフレーションの最も魅力的な特徴は、\textbf{同じポテンシャル $V(\phi)$} から
インフレーションとダークエネルギーの両方が説明できることである。

\begin{keyresult}[同一ポテンシャルからの二重説明]
\begin{itemize}
  \item \textbf{インフレーション期}：$\phi \approx \pi$ の偽真空のエネルギー $V(\pi)$ が
        指数関数的膨張を駆動。
  \item \textbf{現在のダークエネルギー}：$\phi = 0$ の真空エネルギー $V(0) = \delta$ が
        加速膨張を駆動。
  \item 両者は $V(\phi)$ の異なる極小の値に対応。
\end{itemize}
\end{keyresult}

\subsection{Slow-roll パラメータと観測予測}

$\delta$ がダークエネルギースケールで決まるため、インフレーション期の
slow-roll パラメータも $\delta$ に依存する：
\begin{align}
  \epsilon &= \frac{M_{\mathrm{Pl}}^2}{2}\left(\frac{V'}{V}\right)^2, \\
  \eta &= M_{\mathrm{Pl}}^2 \frac{V''}{V}.
\end{align}

これらからスカラースペクトル指数 $n_s$ とテンソル-スカラー比 $r$ が予測される：
\begin{align}
  n_s &= 1 - 6\epsilon + 2\eta, \\
  r &= 16\epsilon.
\end{align}

\begin{warningbox}[$\delta$ 依存性に関する注意]
$n_s$ と $r$ の具体的な数値は $\delta$ の正確な値に強く依存する。
$\delta$ はダークエネルギー密度 $\rho_\Lambda \sim 10^{-47}\,\mathrm{GeV}^4$ と
関係するが、$V(\phi)$ の全体的な形状（kinetic term の正規化を含む）が
まだ厳密に導出されていないため、定量的予測は現時点では不確実である。
これは第18章で未解決問題として議論する。
\end{warningbox}

\begin{hsbox}
インフレーションとは、宇宙が誕生直後（$\sim 10^{-36}$ 秒）に
指数関数的に膨張した時期のこと。これにより、宇宙がなぜこんなに一様で平坦なのかが説明できる。
ダークエネルギーとは、現在の宇宙の加速膨張を引き起こしている未知のエネルギー。
この2つが同じ原因（同じポテンシャル $V(\phi)$）で説明できるなら、
物理学の2大謎が一挙に解決する可能性がある。
\end{hsbox}
  % Ch 11--13: G=-G, gravity dial, domain walls

% ============================================================
%  PART V: 実験プログラム
% ============================================================
\part{実験プログラム}
% ============================================================
%  Part V: 実験プログラム
%  wb_textbook_v2_experiments.tex
%  Chapters 14--16
% ============================================================

% ============================================================
% CHAPTER 14
% ============================================================
\chapter{Phase P --- フォノニック結晶でBerry位相$\pi$を検出}
\label{ch:phase-p}

\begin{experimentbox}[Phase P: Berry位相 $\pi$ の音響検出]
\begin{tabular}{ll}
\textbf{目標} & フォノニック結晶でバルクBerry位相 $\pi$ を音響測定で検出 \\
\textbf{コスト} & 約 {\yen}27,000 \\
\textbf{期間} & 1ヶ月 \\
\textbf{必要設備} & 3Dプリンタ（Bambu Lab P1S推奨）、スピーカー、マイク \\
\textbf{GO条件} & $\pi$ 位相ジャンプの観測
\end{tabular}
\end{experimentbox}

% ------------------------------------------------------------
\section{なぜフォノン系か}
\label{sec:14.1}

第11章で説明したNielsen-Ninomiya定理の抜け穴を再確認する。

電子（フェルミオン）系では、バルク全体のBerry位相を $\pi$ にすることが
Nielsen-Ninomiya定理により困難である。一方、フォノン（ボソン）系には
この制約が適用されない。

フォノニック結晶を使う理由をまとめると：

\begin{enumerate}[label=(\arabic*)]
  \item \textbf{理論的根拠}：Nielsen-Ninomiya定理の制約を受けない（ボソン系）。
  \item \textbf{製造容易性}：3Dプリンタで家庭製作可能。
  \item \textbf{測定容易性}：音響測定は電子測定より遥かに安価で簡単。
  \item \textbf{スケーラビリティ}：ユニットセルサイズ $\sim$ cm、巨視的な試料を作れる。
\end{enumerate}

\begin{hsbox}
フォノニック結晶とは、音に対する「結晶」。
普通の結晶が原子の周期配列で光の振る舞いを制御するように、
フォノニック結晶は構造体の周期配列で音の振る舞いを制御する。
3Dプリンタで作れるサイズ（$\sim$ cm）なので、ガレージ実験が可能。
\end{hsbox}

% ------------------------------------------------------------
\section{Wilson loop 確認済み構造}
\label{sec:14.2}

Berry位相 $\pi$ を実現するフォノニック結晶の設計原理は以下の通りである。

\subsection{2球体ユニットセル}

ユニットセルは\textbf{2つの球体}から構成される。
球体の半径が異なることで\textbf{反転対称性が破れ}、
非自明なBerry位相が生じる。

\begin{dfnbox}[設計パラメータ]
\begin{itemize}
  \item ユニットセルサイズ：$a \approx 20$\,mm（立方体）
  \item 大球半径：$r_1 = 7$\,mm
  \item 小球半径：$r_2 = 4$\,mm
  \item 母材：PLA（3Dプリント用樹脂）
  \item 球体は中空でも中実でもよい
\end{itemize}
\end{dfnbox}

\subsection{Wilson loop 計算}

Wilson loop（ウィルソンループ）は、バンドのBerry位相を計算する標準的手法である。
ブリルアンゾーンの閉じたパスに沿ってBerry接続を積分する：
\[
  \mathcal{W} = \mathcal{P}\exp\left(-i\oint_{\mathcal{C}} \mathbf{A}(\mathbf{k})\cdot d\mathbf{k}\right),
\]
ここで $\mathbf{A}(\mathbf{k})$ はBerry接続、$\mathcal{P}$ は経路順序積である。

この計算を上記の2球体構造に対して実行すると、
最低音響バンドの Wilson loop 固有値が $e^{i\pi} = -1$ となることが確認されている。
すなわち\textbf{Berry位相 $= \pi$}。

\begin{keyresult}[Wilson loop の結果]
2球体フォノニック結晶（反転対称性破れ）のWilson loop計算：
\[
  \phi_{\text{Berry}} = \arg\det\mathcal{W} = \pi.
\]
バルク全体でBerry位相 $\pi$ が実現している。
\end{keyresult}

% ------------------------------------------------------------
\section{3Dプリント手順}
\label{sec:14.3}

\subsection{ワークフロー}

\begin{enumerate}[label=\textbf{Step \arabic*:}]
  \item \textbf{OpenSCAD でモデル作成}：
        2球体ユニットセルを記述するOpenSCADスクリプトを作成。
        $5 \times 5 \times 5 = 125$ ユニットセルの配列を生成。
  \item \textbf{STL エクスポート}：
        OpenSCADからSTLファイルを出力。
        メッシュ品質は $\$fn = 32$ 以上を推奨。
  \item \textbf{スライシング}：
        Bambu Studioでスライス。レイヤー高さ 0.2\,mm。
  \item \textbf{印刷}：
        Bambu Lab P1S（または同等のFDMプリンタ）で印刷。
        PLA使用、サポート材あり。印刷時間 $\sim$ 8時間。
  \item \textbf{後処理}：
        サポート材除去、表面の清掃。
\end{enumerate}

\begin{warningbox}[ジャイロイドインフィルは使用禁止]
スライサーのインフィル（内部充填）パターンに\textbf{ジャイロイドを使ってはならない}。
ジャイロイド構造は反転対称性を保つため、Berry位相が $0$ になってしまう。

\medskip
\textbf{正しいインフィル設定}：インフィル密度 100\%（中実）、
または設計通りの中空構造を使用。スライサーのインフィルパターンに頼らないこと。
\end{warningbox}

% ------------------------------------------------------------
\section{音響透過測定}
\label{sec:14.4}

\subsection{測定セットアップ}

\begin{enumerate}[label=(\alph*)]
  \item \textbf{音源}：スピーカー（周波数掃引 $1$--$20$\,kHz）。
  \item \textbf{検出器}：マイクロフォン（無指向性、周波数応答 $\pm 3$\,dB）。
  \item \textbf{配置}：音源 --- フォノニック結晶 --- 検出器（直線配置）。
  \item \textbf{データ取得}：PCのサウンドカード + Python（scipy.signal）。
\end{enumerate}

\subsection{$\pi$ 位相ジャンプの検出}

バンドギャップの上下でZak位相が $\pi$ だけジャンプする。
これは透過波の\textbf{位相}を測定することで検出できる。

具体的には、バンドギャップ直下の周波数 $f_-$ と直上の周波数 $f_+$ で
透過波の位相差を測定する。Berry位相が $\pi$ の場合：
\[
  \Delta\varphi = \varphi(f_+) - \varphi(f_-) = \pi \quad (\text{mod } 2\pi).
\]

\begin{experimentbox}[Phase P の GO/STOP 判定]
\begin{itemize}
  \item \textbf{GO}：$\Delta\varphi = \pi \pm 0.3$\,rad が観測された場合。
        $\to$ Berry位相 $\pi$ の存在が確認。Phase 0 へ進む。
  \item \textbf{STOP}：$\Delta\varphi \approx 0$ の場合。
        $\to$ 設計の見直し（球体サイズ比、格子定数）。
\end{itemize}
\end{experimentbox}


% ============================================================
% CHAPTER 15
% ============================================================
\chapter{Phase 0--3 --- DNからねじり天秤まで}
\label{ch:phases-0-3}

\begin{experimentbox}[実験フェーズの全体像]
\begin{tabular}{lllr}
\toprule
フェーズ & 検証内容 & 主要装置 & 推定コスト \\
\midrule
Phase 0   & DN/DD基本測定  & ネットワークアナライザ & {\yen}130k \\
Phase 0.5 & $\pi$-接合+BAW統合 & 水晶振動子 & \$10k \\
Phase 1   & 182:1テスト    & BAW共振器 & \$50k \\
Phase 2   & カシミール符号反転 & AFMカンチレバー & \$200k \\
Phase 3   & ねじり天秤重力異常 & Cavendish型天秤 & \$500k \\
\bottomrule
\end{tabular}
\end{experimentbox}

% ------------------------------------------------------------
\section{Phase 0: DN/DD基本測定}
\label{sec:15.1}

Phase 0 はWB理論の最も基本的な予測を検証する。
Dirichlet-Neumann (DN) 境界条件とDirichlet-Dirichlet (DD) 境界条件の
スペクトルの差異を音響共振器で測定する。

\begin{itemize}
  \item \textbf{装置}：アルミニウム管（長さ $L \sim 1$\,m）+ マイク + スピーカー。
  \item \textbf{DN条件}：一端開放、他端閉鎖。
  \item \textbf{DD条件}：両端閉鎖。
  \item \textbf{測定量}：共鳴周波数列 $f_n$ のパターン。
  \item \textbf{推定コスト}：約 {\yen}130,000（ネットワークアナライザ含む）。
\end{itemize}

DN条件とDD条件では共鳴周波数のパターンが異なり、
そのスペクトルの差がゼータ関数の零点構造を反映するかを確認する。

% ------------------------------------------------------------
\section{Phase 0.5: $\pi$-接合 + BAW統合}
\label{sec:15.2}

Phase Pで確認した $\pi$-Berry位相を、圧電素子（BAW: Bulk Acoustic Wave）の
共振器と統合する。

\begin{itemize}
  \item \textbf{装置}：水晶振動子（AT-cut quartz）+ $\pi$-Berry結晶の接合体。
  \item \textbf{目標}：$\pi$-Berry結晶が水晶振動子の共鳴周波数にシフトを与えるかを測定。
  \item \textbf{予測}：Berry位相 $\pi$ による境界条件変化が、
        共鳴周波数を $\Delta f / f \sim 10^{-6}$ 程度シフトさせる。
  \item \textbf{推定コスト}：約 \$10,000。
\end{itemize}

% ------------------------------------------------------------
\section{Phase 1: 182:1 テスト}
\label{sec:15.3}

\begin{keyresult}[$K = \zeta$ の実験検証]
Phase 1 は WB 理論の核心的予測を検証する。
BC真空の分配関数 $K(t)$ がリーマンゼータ関数 $\zeta(t)$ と一致するならば、
DN/DD スペクトル比が特定の値 $182:1$ に近い値を取るはずである。
これは公理2の直接検証である。
\end{keyresult}

\begin{itemize}
  \item \textbf{装置}：高精度BAW共振器（$Q > 10^6$）。
  \item \textbf{測定}：DN境界とDD境界のスペクトルを高精度で比較。
  \item \textbf{予測}：$\zeta$ 関数の Euler 積表現が物理的に成立するなら、
        スペクトルの素数構造が $182:1$ の比として現れる。
  \item \textbf{推定コスト}：約 \$50,000。
  \item \textbf{意義}：成功すれば、数論（Euler積）が物理現象として実在する初の証拠。
\end{itemize}

\begin{hsbox}
$182 = 2 \times 7 \times 13$。この数は $\zeta$ 関数の特殊値から計算で出てくるが、
この数値そのものよりも「計算値と実測値が一致する」ことが重要。
1/3 とのずれがレギュレータ効果で説明できるかが鍵になる（仮説 H5）。
\end{hsbox}

% ------------------------------------------------------------
\section{Phase 2: $\pi$-接合カシミール符号反転}
\label{sec:15.4}

カシミール効果は、2枚の平行な導体板の間に量子真空のゆらぎが生む力である。
通常は\textbf{引力}だが、Berry位相 $\pi$ の物質を使うと\textbf{斥力}に転じると予測される。

\begin{itemize}
  \item \textbf{装置}：AFM（原子間力顕微鏡）カンチレバー型カシミール力測定装置。
  \item \textbf{試料}：一方の板をトポロジカル絶縁体（Bi$_2$Se$_3$等）に交換。
  \item \textbf{予測}：カシミール力の符号反転（引力 $\to$ 斥力）。
  \item \textbf{推定コスト}：約 \$200,000（AFM装置 + TI試料）。
  \item \textbf{意義}：Berry位相が真空のゼロ点エネルギーに影響することの直接証拠。
\end{itemize}

\begin{warningbox}[技術的困難]
カシミール力の精密測定は極めて難しい。表面粗さ、静電パッチ効果、温度効果など
多くの系統誤差源がある。符号反転の確認には、通常金属板との比較測定が不可欠。
\end{warningbox}

% ------------------------------------------------------------
\section{Phase 3: ねじり天秤重力異常}
\label{sec:15.5}

最終段階。$\pi$-Berry物質の近傍で、ニュートン重力定数 $G$ の
異常（$\Geff \neq G$）をねじり天秤で直接検出する。

\begin{itemize}
  \item \textbf{装置}：Cavendish型ねじり天秤（感度 $\Delta G/G \sim 10^{-4}$）。
  \item \textbf{試料}：$\pi$-Berry物質のバルク試料（$\sim$ cm$^3$）。
  \item \textbf{測定}：$\pi$-Berry試料 vs 通常試料での重力力の比較。
  \item \textbf{予測}：$\Geff = -G$（完全反転）。ただし、Berry位相が物質表面で
        急速に $0$ に緩和する場合、効果は大幅に減衰する可能性。
  \item \textbf{推定コスト}：約 \$500,000。
  \item \textbf{意義}：$\Geff = -G$ の\textbf{直接検証}。成功すれば物理学のパラダイムシフト。
\end{itemize}

% ------------------------------------------------------------
\section{GO/STOP判定フロー}
\label{sec:15.6}

各フェーズは前段階の成功を条件として進行する。

\begin{center}
\begin{tikzpicture}[
  node distance=1.5cm,
  box/.style={draw, rounded corners, minimum width=3cm, minimum height=0.8cm,
              fill=blue!10, font=\small},
  go/.style={-{Stealth}, thick, green!60!black},
  stop/.style={-{Stealth}, thick, red!60!black},
]
  \node[box] (PP) {Phase P};
  \node[box, right=2.5cm of PP] (P0) {Phase 0};
  \node[box, right=2.5cm of P0] (P05) {Phase 0.5};
  \node[box, below=1.2cm of PP] (P1) {Phase 1};
  \node[box, right=2.5cm of P1] (P2) {Phase 2};
  \node[box, right=2.5cm of P2] (P3) {Phase 3};

  \draw[go] (PP) -- node[above,font=\scriptsize]{GO} (P0);
  \draw[go] (P0) -- node[above,font=\scriptsize]{GO} (P05);
  \draw[go] (P05) -- node[right,font=\scriptsize]{GO} (P1);
  \draw[go] (P1) -- node[above,font=\scriptsize]{GO} (P2);
  \draw[go] (P2) -- node[above,font=\scriptsize]{GO} (P3);

  \node[below=0.3cm of PP, red!60!black, font=\scriptsize] {STOP: 再設計};
  \node[below=0.3cm of P1, red!60!black, font=\scriptsize] {STOP: H5見直し};
  \node[below=0.3cm of P2, red!60!black, font=\scriptsize] {STOP: H4見直し};
\end{tikzpicture}
\end{center}

\begin{gostopbox}[判定基準のまとめ]
\begin{itemize}
  \item \textbf{Phase P STOP} $\to$ フォノニック結晶の再設計。理論自体は否定されない。
  \item \textbf{Phase 1 STOP} $\to$ 仮説H5（レギュレータ予想）の見直し。$K = \zeta$ 自体は生存可能。
  \item \textbf{Phase 2 STOP} $\to$ 仮説H4（Berry位相$\to$重力結合）の否定。5定数の導出は生存。
  \item \textbf{Phase 3 STOP} $\to$ 仮説H3（スペクトル重力）の否定、または効果が検出限界以下。
\end{itemize}
\end{gostopbox}


% ============================================================
% CHAPTER 16
% ============================================================
\chapter{何を検証しているのか}
\label{ch:what-we-test}

\begin{keyresult}[本章の主題]
各実験フェーズが理論体系のどの部分を検証しているかを明確にする。
実験の成功/失敗が何を意味し、何を意味しないかの論理的整理。
\end{keyresult}

各実験は WB 理論の異なる「層」を検証する。ある実験が失敗しても、
他の予測が無傷で生き残る場合がある。この構造を正確に理解することが重要である。

% ------------------------------------------------------------
\section{Phase P が検証するもの}
\label{sec:16.1}

\begin{experimentbox}[Phase P $\to$ Berry位相 $\pi$ の存在]
\textbf{検証対象}：フォノニック結晶でバルクBerry位相 $\pi$ が実在すること。

\medskip
\textbf{WB公理との関係}：3公理とは\textbf{独立}。Berry位相 $\pi$ の存在は
凝縮系物理学の範囲内の問題であり、算術的時空仮説の検証ではない。

\medskip
\textbf{成功の意味}：Nielsen-Ninomiya定理のボソン系抜け穴の実証。
後続フェーズの\textbf{必要条件}が満たされる。

\medskip
\textbf{失敗の意味}：設計の問題であり、理論の否定ではない。再設計して再試行。
\end{experimentbox}

% ------------------------------------------------------------
\section{Phase 1 が検証するもの}
\label{sec:16.2}

\begin{experimentbox}[Phase 1 $\to$ Euler積の物理性（公理2の検証）]
\textbf{検証対象}：BC真空の分配関数 $K(t)$ がリーマンゼータ関数 $\zeta(t)$ と
一致するかどうか。具体的には、$182:1$ 比の検出。

\medskip
\textbf{WB公理との関係}：\textbf{公理2}（素数独立性）の直接検証。
Euler積 $\zeta(s) = \prod_p (1-p^{-s})^{-1}$ が物理的スペクトルに反映されるかを問う。

\medskip
\textbf{成功の意味}：数論的構造（Euler積）が物理現象として実在する初の証拠。
WB理論の数学的基盤が物理的に正しいことの強力な傍証。

\medskip
\textbf{失敗の意味}：仮説H5（レギュレータ予想：182 vs 1/3）の否定。
ただし $K = \zeta$ 自体が否定されるわけではない（レギュレータの形が違うだけかもしれない）。
\end{experimentbox}

% ------------------------------------------------------------
\section{Phase 2 が検証するもの}
\label{sec:16.3}

\begin{experimentbox}[Phase 2 $\to$ Berry位相が真空に影響（スペクトル作用の検証）]
\textbf{検証対象}：Berry位相 $\pi$ がカシミール効果（真空のゼロ点エネルギー）を
変調するかどうか。

\medskip
\textbf{WB公理との関係}：スペクトル作用原理（仮説H3）の検証。
スペクトル三つ組の自己同型が物理的真空に作用するかを問う。

\medskip
\textbf{成功の意味}：非可換幾何学のスペクトル作用が実験的に検証される初の事例。
反重力理論への決定的一歩。

\medskip
\textbf{失敗の意味}：仮説H4（Berry位相 $\to$ 重力結合）の否定の可能性。
ただし、5定数の導出（$\Omega_\Lambda$ 等）はBerry位相とは独立なので\textbf{生存}する。
\end{experimentbox}

% ------------------------------------------------------------
\section{Phase 3 が検証するもの}
\label{sec:16.4}

\begin{experimentbox}[Phase 3 $\to$ $\Geff = -G$（反重力の直接検証）]
\textbf{検証対象}：$\pi$-Berry物質の近傍でニュートン重力定数の符号反転が起こるか。

\medskip
\textbf{WB公理との関係}：仮説H3（スペクトル重力）+ 仮説H4（Berry $\to$ 重力）の
\textbf{組み合わせ}検証。両方が正しい場合のみ $\Geff = -G$ が実現。

\medskip
\textbf{成功の意味}：反重力の実証。物理学のパラダイムシフト。

\medskip
\textbf{失敗の意味}：H3またはH4（またはその両方）の否定。
ただし $\Omega_\Lambda = 2\pi/9$ や $\alpha_{\mathrm{EM}}$ の導出は独立に生存する。
\end{experimentbox}

% ------------------------------------------------------------
\section{Euclid 2028 が検証するもの}
\label{sec:16.5}

\begin{experimentbox}[Euclid 2028 $\to$ $\Omega_\Lambda = 2\pi/9$（公理3の検証）]
\textbf{検証対象}：ダークエネルギー密度パラメータ $\Omega_\Lambda$ が
$2\pi/9 \approx 0.6981$ であるかどうか。

\medskip
\textbf{WB公理との関係}：\textbf{公理3}（民主主義原理）の検証。
$\Omega_\Lambda = 2\pi/9$ は3公理から直接導出される。

\medskip
\textbf{検証方法}：ESA Euclid衛星（2023年打ち上げ済み）のデータリリース（2028年予定）。
弱重力レンズ + バリオン音響振動の組み合わせで $\Omega_\Lambda$ を $\pm 0.005$ の精度で決定予定。

\medskip
\textbf{成功の意味}：算術的時空仮説の宇宙論的予測が正しいことの証拠。
$2\sigma$ 以内に $2\pi/9$ が含まれれば強力な支持。

\medskip
\textbf{失敗の意味}：公理3（民主主義）の否定、またはH1c（民主主義仮説の細部）の修正が必要。
ただし、他の4定数の導出は $\Omega_\Lambda$ とは独立な論理で成り立つ部分もある。
\end{experimentbox}

\begin{hsbox}
Euclid（ユークリッド）はESA（欧州宇宙機関）の宇宙望遠鏡。
宇宙の大規模構造を精密に測定し、ダークエネルギーの性質を明らかにすることが目標。
WB理論は「$\Omega_\Lambda$ はちょうど $2\pi/9$」と予測しており、
これは現在の観測値 $0.685 \pm 0.013$ と $1\sigma$ 以内で一致している。
Euclid のデータでこの一致が確認されるかが、2028年の最大の検証ポイントである。
\end{hsbox}
  % Ch 14--16: Phase P, Phase 0--3

% ============================================================
%  PART VI: 誠実な評価と未解決問題
% ============================================================
\part{誠実な評価と未解決問題}
% ============================================================
%  Part VI: 誠実な評価と未解決問題
%  wb_textbook_v2_honest.tex
%  Chapters 17--18
% ============================================================

% ============================================================
% CHAPTER 17
% ============================================================
\chapter{仮説依存ネットワーク}
\label{ch:dependency}

\begin{keyresult}[本章の主題]
MA理論は複数の仮説の\textbf{積み重ね}で構成されている。
どの仮説が崩れると何が倒れるかを、依存関係図として明示する。
科学理論の信頼性は、その透明性に比例する。
\end{keyresult}

% ------------------------------------------------------------
\section{レベル0：数学的事実}
\label{sec:17.1}

以下は\textbf{証明済みの数学的定理}であり、物理的仮説とは独立に成立する。
これらが崩れることはない。

\begin{dfnbox}[数学的事実 M1--M6]
\begin{description}[style=nextline, leftmargin=1.5cm]
  \item[M1] リーマンゼータ関数の関数等式：
    $\zeta(s) = 2^s \pi^{s-1}\sin(\pi s/2)\,\Gamma(1-s)\,\zeta(1-s)$。
  \item[M2] オイラー積公式：
    $\zeta(s) = \prod_p (1-p^{-s})^{-1}$（$\Re(s) > 1$）。
  \item[M3] Connes の BC 系の分類定理：
    BC 代数の KMS 状態は臨界温度 $\beta = 1$ で相転移を示す。
  \item[M4] 非可換幾何学のスペクトル作用原理：
    スペクトル三つ組 $(\mathcal{A}, \mathcal{H}, D)$ からボゾン作用が復元される。
  \item[M5] Nielsen-Ninomiya 定理：
    格子フェルミオンにはダブラーが存在する（ボソン系には不適用）。
  \item[M6] Berry位相の量子化：
    反転対称性の破れた系で Zak 位相は $0$ または $\pi$。
\end{description}
\end{dfnbox}

% ------------------------------------------------------------
\section{レベル1：物理的仮説}
\label{sec:17.2}

以下は\textbf{未証明の物理的仮説}である。実験で検証（または反証）される対象。

\subsection{H1: 算術時空仮説}

「時空は $\Spec(\Z)$ 上のアデール幾何として記述される」という中心仮説。
これはさらに3つの下位仮説に分解される。

\begin{hypbox}[H1: 算術時空（3公理に分解）]
\begin{description}[style=nextline, leftmargin=2cm]
  \item[H1a（離散性）] 時空の基本構造は離散的（$\Spec(\Z)$）であり、連続時空は創発的。
  \item[H1b（素数独立）] 各素数 $p$ は\textbf{独立な}局所因子 $\Q_p$ を与える。
        物理は全ての素数からの寄与の「積」（Euler積構造）として表される。
  \item[H1c（民主主義）] 全ての素数は\textbf{等しい重み}で寄与する。
        特別な素数は存在しない。
\end{description}
\end{hypbox}

\subsection{H2: BC真空}

\begin{hypbox}[H2: BC真空（$K = \zeta$）]
Bost-Connes系の分配関数 $K(t)$ が宇宙の真空のスペクトルを記述する：
\[
  K(t) = \zeta(t) = \sum_{n=1}^{\infty} n^{-t} = \prod_p (1 - p^{-t})^{-1}.
\]
これは H1 の3公理から\textbf{導出される}（公理 $\to$ BC系 $\to$ $K = \zeta$）。
\end{hypbox}

\subsection{H3: スペクトル重力}

\begin{hypbox}[H3: スペクトル重力]
非可換幾何学のスペクトル作用原理が物理的に正しい：
スペクトル三つ組 $(\mathcal{A}, \mathcal{H}, D)$ から Einstein 方程式が導出される。
これは Connes-Chamseddine の定理（M4）の\textbf{物理的適用}を仮説とする。
\end{hypbox}

\subsection{H4: Berry位相 $\to$ 重力結合}

\begin{hypbox}[H4: Berry位相と重力の結合]
物質中のBerry位相 $\phi$ がスペクトル三つ組の自己同型を誘導し、
有効重力定数を変調する：
\[
  f_k \;\to\; e^{-i\phi(d-k)/2}\, f_k
  \qquad\Longrightarrow\qquad
  \Geff = G\cos\phi.
\]
\end{hypbox}

\subsection{H5: レギュレータ予想}

\begin{hypbox}[H5: レギュレータ予想（182 vs 1/3）]
BC系のスペクトルから導出されるDN/DDスペクトル比は $182:1$ であるが、
物理的に観測される比は正則化（レギュレータ）効果により $1/3$ に近い値となる。
この正則化の機構と数値が正しいかは未検証。
\end{hypbox}

% ------------------------------------------------------------
\section{依存関係図}
\label{sec:17.3}

以下の図は、仮説間の論理的依存関係を示す。矢印は「$A \to B$」で
「$A$ が成立しないと $B$ は（少なくとも現在の論証では）成立しない」を意味する。

\begin{center}
\begin{tikzpicture}[
  node distance=1.8cm and 2.5cm,
  hyp/.style={draw, rounded corners, minimum width=2.2cm, minimum height=0.7cm,
              fill=purple!10, font=\small\bfseries},
  math/.style={draw, rounded corners, minimum width=2.2cm, minimum height=0.7cm,
               fill=green!10, font=\small},
  pred/.style={draw, rounded corners, minimum width=2.8cm, minimum height=0.7cm,
               fill=yellow!15, font=\small},
  arr/.style={-{Stealth}, thick},
]
  % Math level
  \node[math] (M3) {M3: BC分類};
  \node[math, left=1.5cm of M3] (M2) {M2: Euler積};
  \node[math, right=1.5cm of M3] (M4) {M4: スペクトル作用};
  \node[math, right=1.5cm of M4] (M6) {M6: Berry量子化};

  % Hypothesis level
  \node[hyp, below=1.5cm of M2] (H1) {H1: 算術時空};
  \node[hyp, below=1.5cm of M3] (H2) {H2: BC真空};
  \node[hyp, below=1.5cm of M4] (H3) {H3: スペクトル重力};
  \node[hyp, below=1.5cm of M6] (H4) {H4: Berry$\to$重力};
  \node[hyp, below=0.8cm of H2] (H5) {H5: レギュレータ};

  % Prediction level
  \node[pred, below=2.5cm of H1] (OL) {$\Omega_\Lambda = 2\pi/9$};
  \node[pred, right=0.5cm of OL] (aEM) {$\alpha_{\mathrm{EM}}$};
  \node[pred, below=2.5cm of H3] (Geff) {$\Geff = -G$};
  \node[pred, below=2.5cm of H4] (Cas) {カシミール反転};
  \node[pred, below=1.5cm of H5] (r182) {182:1比};

  % Arrows
  \draw[arr] (M2) -- (H1);
  \draw[arr] (M3) -- (H2);
  \draw[arr] (H1) -- (H2);
  \draw[arr] (M4) -- (H3);
  \draw[arr] (M6) -- (H4);
  \draw[arr] (H2) -- (H5);

  \draw[arr] (H1) -- (OL);
  \draw[arr] (H2) -- (OL);
  \draw[arr] (H2) -- (aEM);
  \draw[arr] (H3) -- (Geff);
  \draw[arr] (H4) -- (Geff);
  \draw[arr] (H4) -- (Cas);
  \draw[arr] (H3) -- (Cas);
  \draw[arr] (H5) -- (r182);

\end{tikzpicture}
\end{center}

% ------------------------------------------------------------
\section{否定シナリオ：何が崩れると何が倒れるか}
\label{sec:17.4}

仮説の否定が波及する範囲を整理する。

\subsection{H3 が否定された場合}

\begin{warningbox}[H3否定の波及]
\textbf{スペクトル重力が物理的に成立しない場合}：
\begin{itemize}
  \item \textbf{全滅}：$\Geff = -G$（反重力）、カシミール符号反転、Berry位相インフレーション。
        $\to$ 反重力関連の予測は\textbf{全て崩壊}する。
  \item \textbf{生存}：$\Omega_\Lambda = 2\pi/9$、$\alpha_{\mathrm{EM}}$ の導出、$m_p/m_e$ の導出。
        $\to$ これらは H1 + H2 から導出されており、H3 とは\textbf{独立}。
\end{itemize}
\textbf{要点}：反重力は死ぬが、宇宙定数の予測は生き残る。
\end{warningbox}

\subsection{H4 が否定された場合}

\begin{warningbox}[H4否定の波及]
\textbf{Berry位相が重力と結合しない場合}：
\begin{itemize}
  \item \textbf{全滅}：$\Geff = -G$（反重力）、重力ダイヤル、カシミール符号反転。
        $\to$ 反重力の\textbf{工学的応用}は全て不可能。
  \item \textbf{生存}：5定数の導出（$\Omega_\Lambda$, $\alpha_{\mathrm{EM}}$, $m_p/m_e$, $\theta_W$, $\delta_{\mathrm{CP}}$）。
        $\to$ 定数の導出は Berry 位相工学とは\textbf{独立}な論理に基づく。
\end{itemize}
\textbf{要点}：重力制御の夢は消えるが、「なぜ物理定数がこの値か」への解答は生き残る。
\end{warningbox}

\subsection{H1c が否定された場合}

\begin{warningbox}[H1c否定の波及]
\textbf{素数の民主主義が成立しない場合}（一部の素数が特別な重みを持つ）：
\begin{itemize}
  \item \textbf{影響}：$\Omega_\Lambda = 2\pi/9$ の値が\textbf{ずれる}。
        正確な値は重みの分布に依存。
  \item \textbf{生存}：$\Omega_\Lambda$ が「何らかの数論的値」であることは維持される。
        5定数間の\textbf{関係式}も多くが生存。
  \item \textbf{判定}：Euclid 2028 の $\Omega_\Lambda$ 精密測定で決着。
        $2\pi/9$ から有意にずれれば H1c は否定される。
\end{itemize}
\end{warningbox}

\begin{hsbox}
科学理論の美しさは「反証可能であること」にある。
MA理論は多くの仮説の積み重ねだが、各仮説が\textbf{独立に}検証・反証可能な構造になっている。
全てが正しい必要はなく、一部が崩れても他の部分が生き残る。
これは「トランプの家」ではなく「レンガの壁」に近い構造である。
\end{hsbox}


% ============================================================
% CHAPTER 18
% ============================================================
\chapter{失敗の記録と未解決問題}
\label{ch:failures}

\begin{warningbox}[本章の精神]
科学では成功だけでなく\textbf{失敗を記録する}ことが不可欠である。
失敗を隠蔽すれば、他の研究者が同じ誤りを繰り返す。
本章はMAプログラムの失敗と未解決問題を正直に記録する。
\end{warningbox}

% ------------------------------------------------------------
\section{失敗リスト}
\label{sec:18.1}

以下は、研究過程で発生した具体的な誤りと失敗を時系列で記録したものである。

\subsection{失敗1：符号誤り（3D DN = 負エネルギーではなかった）}

\begin{warningbox}[失敗1: 3D DN の符号]
\textbf{誤り}：3次元のDirichlet-Neumann境界条件が負のカシミールエネルギーを
与えると考え、これを反重力の証拠として議論した。

\medskip
\textbf{実際}：3D DN のカシミールエネルギーの符号は正（通常の引力）であった。
符号が反転するのは特定の次元・幾何のみであり、3D の一般的な場合には当てはまらない。

\medskip
\textbf{教訓}：カシミール効果の符号は幾何と次元に強く依存する。
安易な一般化は危険。
\end{warningbox}

\subsection{失敗2：$10^{43}$\,J 見積もり誤り}

\begin{warningbox}[失敗2: 壁エネルギーの過大評価]
\textbf{誤り}：ドメインウォールのエネルギーを $\sim 10^{43}$\,J と見積もり、
実用的な反重力装置には天文学的エネルギーが必要と結論した。

\medskip
\textbf{実際}：正しい見積もりは $\sim 10^3$\,J（キロジュール程度）。
プランクスケールの壁ではなく、ダークエネルギースケールの壁であることを
見落としていた（40桁の過大評価！）。

\medskip
\textbf{教訓}：エネルギースケールの同定を慎重に行うこと。
\end{warningbox}

\subsection{失敗3：$G = -2G$（Liouville）は無限操作が必要}

\begin{warningbox}[失敗3: Liouville モード]
\textbf{誤り}：Liouville モードの共形因子を利用して $\Geff = -2G$ を
実現できると考えた。

\medskip
\textbf{実際}：この操作には\textbf{無限回}の共形変換の合成が必要であり、
物理的に実現不可能。有限回の操作では $\Geff = -2G$ は達成できない。

\medskip
\textbf{教訓}：数学的に形式的に成り立つ操作が物理的に実現可能とは限らない。
\end{warningbox}

\subsection{失敗4：ジャイロイドインフィルは Berry 位相 0}

\begin{warningbox}[失敗4: ジャイロイド構造]
\textbf{誤り}：3Dプリントのジャイロイドインフィルパターンがトポロジカル
フォノニック結晶として機能すると期待した。

\medskip
\textbf{実際}：ジャイロイド構造は空間反転対称性を保存するため、
Zak 位相（Berry 位相）は\textbf{厳密に $0$}。トポロジカルに自明。

\medskip
\textbf{教訓}：対称性の解析を設計前に行うこと。
反転対称性を\textbf{破る}構造が必須。
\end{warningbox}

\subsection{失敗5：DESI $w_0 w_a$ フィット失敗}

\begin{warningbox}[失敗5: 状態方程式のフィット]
\textbf{誤り}：DESI（Dark Energy Spectroscopic Instrument）の $w_0$-$w_a$
パラメータ化を、Berry 位相ポテンシャル $V(\phi)$ からフィットしようとした。

\medskip
\textbf{実際}：検討した\textbf{全てのポテンシャル形状}で、$w_a$ の傾斜が
DESI データと合わない。$w_0 \approx -1$ は再現できるが、
$w_a$ の非ゼロ成分（時間変化）は正しく再現できなかった。

\medskip
\textbf{教訓}：$\Lambda$CDM（$w_0 = -1$, $w_a = 0$）からのずれは、
Berry 位相インフレーションの単純な版では説明困難。
\end{warningbox}

\subsection{失敗6：$\alpha_s$ の導出失敗}

\begin{warningbox}[失敗6: 強い結合定数]
\textbf{誤り}：電磁結合定数 $\alpha_{\mathrm{EM}}$ と同様の手法で
強い結合定数 $\alpha_s$ を導出しようとした。

\medskip
\textbf{実際}：最良の試みでも実験値から\textbf{4.5\%のずれ}が残る。
$\alpha_{\mathrm{EM}}$ の $0.3\%$ 精度に比べて大幅に悪い。

\medskip
\textbf{教訓}：$\alpha_s$ のランニング（エネルギースケール依存性）と
非摂動論的効果が、単純な数論的導出を困難にしている可能性。
\end{warningbox}

\subsection{失敗7：$m_\mu / m_e$ の 23 が未導出}

\begin{warningbox}[失敗7: ミュー粒子質量比の因子 23]
\textbf{誤り}：$m_\mu / m_e \approx 206.768$ を完全に導出しようとした。

\medskip
\textbf{実際}：$m_\mu / m_e$ の値に現れる因子 $23$ の数論的起源が未解明。
他の因子は $\zeta$ 関数の特殊値や $\pi$ から説明できるが、
$23$ だけが「浮いて」いる。

\medskip
\textbf{教訓}：レプトン質量階層の完全な理解には、まだ欠けている要素がある。
\end{warningbox}

\subsection{失敗8：浮遊する球のSTL（印刷不可能）}

\begin{warningbox}[失敗8: 印刷不可能な設計]
\textbf{誤り}：フォノニック結晶の初期設計で、
ユニットセル内の球が母材と接続されていない「浮遊する球」を含むSTLファイルを生成。

\medskip
\textbf{実際}：FDM 3Dプリンタでは、浮遊する構造は\textbf{印刷不可能}
（サポート材を除去すると球が落下する）。
球を細い柱で母材に接続する設計変更が必要だった。

\medskip
\textbf{教訓}：計算物理と製造工学の間にはギャップがある。
印刷可能性（manufacturability）を設計段階で検討すること。
\end{warningbox}

% ------------------------------------------------------------
\section{未解決問題}
\label{sec:18.2}

以下は現時点で解決されていない理論的課題である。

\subsection{未解決問題1：$V(\phi)$ の厳密導出}

\begin{dfnbox}[未解決: ポテンシャルの導出]
Berry位相スカラー・テンソル理論のポテンシャル $V(\phi)$ を
スペクトル三つ組 $(\mathcal{A}, \mathcal{H}, D)$ から\textbf{厳密に}導出すること。

\medskip
\textbf{現状}：$V(\phi)$ の定性的な形（2つの極小）は分かっているが、
kinetic term の正規化を含めた定量的な形状が決まっていない。
これにより、slow-roll パラメータ（$n_s$, $r$）の定量的予測ができない。

\medskip
\textbf{難しさ}：スペクトル作用の展開で、高次項の寄与が kinetic term の
係数を修正する。この修正の正確な値は非可換幾何学の未解決問題と絡む。
\end{dfnbox}

\subsection{未解決問題2：$\alpha_s$}

\begin{dfnbox}[未解決: 強い結合定数]
QCD の結合定数 $\alpha_s(M_Z) \approx 0.1179$ を数論的に導出すること。

\medskip
\textbf{現状}：最良の試みで $4.5\%$ のずれ（失敗6参照）。
$\alpha_{\mathrm{EM}}$ の $0.3\%$ 精度に遠く及ばない。

\medskip
\textbf{可能性}：$\alpha_s$ の大きさ（$\sim 0.1$、非摂動論的領域に近い）が
単純な摂動論的数論的手法の適用を困難にしている可能性がある。
非摂動論的効果（instanton、confinement）を数論的に取り込む必要があるかもしれない。
\end{dfnbox}

\subsection{未解決問題3：$\alpha_{\mathrm{EM}}$ の $0.3\%$ ずれ}

\begin{dfnbox}[未解決: 電磁結合定数の微調整]
導出値 $\alpha_{\mathrm{EM}}^{-1} \approx 137.02$ と実験値 $137.036$ の
$0.3\%$ のずれの起源。

\medskip
\textbf{仮説}：ランニング効果（繰り込み群フロー）による補正。
低エネルギーでの値は高エネルギーの「裸の値」からランニングで決まるが、
MA理論が与えるのはどのスケールの値かが不明確。

\medskip
\textbf{検証方法}：ランニングの計算を行い、MA理論の値がどのエネルギースケールに
対応するかを同定する。$\mu \sim M_{\mathrm{Pl}}$ での値なら、
$\mu \sim m_e$ への繰り込み群フローで $0.3\%$ のずれが説明できるかもしれない。
\end{dfnbox}

\subsection{未解決問題4：$d = 4$ の導出}

\begin{dfnbox}[未解決: なぜ4次元か]
時空がなぜ4次元（$3+1$次元）であるかの導出。

\medskip
\textbf{現状}：MA理論は $d = 4$ を\textbf{仮定}しており、導出していない。
$\Spec(\Z)$ は1次元であり、ここから $d = 4$ の時空が創発する機構は不明。

\medskip
\textbf{手がかり}：
\begin{itemize}
  \item $\Spec(\Z)$ のコホモロジー次元 $\cd(\Spec(\Z)) = 3$（Arakelov補正込み）
        $\to$ これが $d = 3+1$ と関係する可能性。
  \item 非可換幾何学の KO 次元が $d = 4 \pmod{8}$ を好む
        $\to$ フェルミオンの自由度の制約。
  \item 人間原理的議論：$d \neq 4$ では安定な原子・惑星軌道が存在しない。
\end{itemize}

\medskip
\textbf{率直な評価}：これは数学と物理学の最も深い未解決問題の一つであり、
MA理論の枠内で短期間に解決できる見込みは低い。
しかし、$\Spec(\Z)$ のコホモロジー構造が鍵を握る可能性はある。
\end{dfnbox}

\begin{hsbox}
「なぜ3次元空間なのか」は子供でも思いつく問いだが、
物理学者にとっても答えのない根本問題である。
前後左右上下の3方向＋時間の1方向で $3+1 = 4$ 次元。
もし空間が2次元なら消化管が体を分断してしまうし、
5次元以上なら惑星軌道が不安定になる。
しかし「だから4次元」というのは答えではなく、結果の言い換えにすぎない。
\end{hsbox}

% ------------------------------------------------------------
\section{まとめ：誠実さの価値}
\label{sec:18.3}

本章で列挙した8つの失敗と4つの未解決問題は、MA理論の\textbf{弱点}である。
しかし弱点を公開することは弱さではなく、科学的\textbf{強さ}の証明である。

\begin{keyresult}[MA理論の現在地]
\begin{itemize}
  \item \textbf{成功}：5つの物理定数の導出（$\Omega_\Lambda$, $\alpha_{\mathrm{EM}}$, $m_p/m_e$, $\theta_W$, $\delta_{\mathrm{CP}}$）、
        Berry位相による $\Geff = G\cos\phi$ の定式化、
        実験可能な予測の体系。
  \item \textbf{失敗}：8件の具体的誤りを修正済み。
  \item \textbf{未解決}：$V(\phi)$ の厳密導出、$\alpha_s$、$\alpha_{\mathrm{EM}}$ の微調整、$d = 4$ の起源。
  \item \textbf{次のステップ}：Phase P 実験（{\yen}27,000、1ヶ月）で
        Berry位相 $\pi$ の音響検出から開始。
\end{itemize}

これは完成された理論ではなく、\textbf{進行中の研究プログラム}である。
読者には批判的な目で本書を読み、誤りを見つけたら報告してほしい。
科学は協力によって進む。
\end{keyresult}
       % Ch 17--18: dependency network, failures, open problems

\end{document}
